\documentclass{book}
\usepackage{amsmath,amssymb}

\begin{document}

\tableofcontents

\chapter{Sistemi di equazioni lineari}
\section{Un primo problema economico: domanda, offerta ed equilibrio}\label{sec:domanda-offerta}

\subsection{Il problema economico}

Uno dei problemi più semplici dell'analisi economica consiste nella
determinazione del prezzo e della quantità di equilibrio di un mercato.
Nonostante la sua semplicità, questo problema costituisce un utile punto
di partenza per introdurre alcuni degli strumenti computazionali che
utilizzeremo nel corso.

Consideriamo un mercato nel quale la domanda e l'offerta sono descritte
dalle seguenti funzioni lineari:

\begin{equation}
    Q_d = \bar{Q}_d - aP,
    \label{eq:domanda}
\end{equation}

e

\begin{equation}
    Q_s = \bar{Q}_s + bP,
    \label{eq:offerta}
\end{equation}

dove $Q_d$ rappresenta la quantità domandata, $Q_s$ la quantità
offerta e $P$ il prezzo del bene.

I parametri $\bar{Q}_d$ e $\bar{Q}_s$ rappresentano rispettivamente
le intercette delle funzioni di domanda e di offerta, mentre $a$ e
$b$ ne determinano la sensibilità rispetto al prezzo.

Assumiamo

\begin{equation}
    a > 0,
    \qquad
    b > 0.
\end{equation}

La quantità domandata diminuisce quindi all'aumentare del prezzo,
mentre la quantità offerta aumenta.

L'equilibrio del mercato si verifica quando la quantità che i
consumatori desiderano acquistare coincide con quella che i produttori
desiderano vendere. Indicando la quantità di equilibrio con $Q$, la
condizione di equilibrio è pertanto

\begin{equation}
    Q_d = Q_s = Q.
    \label{eq:equilibrio}
\end{equation}

Le due variabili endogene del modello sono dunque $Q$ e $P$.
Una volta assegnati i valori dei parametri
$\bar{Q}_d$, $\bar{Q}_s$, $a$ e $b$, il nostro problema consiste nel
determinare il prezzo e la quantità che soddisfano simultaneamente
le equazioni di domanda e di offerta.


\subsection{Un esempio numerico}

Consideriamo la seguente parametrizzazione:

\begin{equation}
    \bar{Q}_d = 1000,
    \qquad
    \bar{Q}_s = 250,
    \qquad
    a = 10,
    \qquad
    b = 5.
\end{equation}

Le funzioni di domanda e di offerta diventano quindi

\begin{equation}
    Q_d = 1000 - 10P
\end{equation}

e

\begin{equation}
    Q_s = 250 + 5P.
\end{equation}

Per determinare l'equilibrio imponiamo l'uguaglianza tra domanda
e offerta:

\begin{equation}
    1000 - 10P = 250 + 5P.
\end{equation}

Riorganizzando i termini otteniamo

\begin{equation}
    750 = 15P,
\end{equation}

da cui segue

\begin{equation}
    P^* = 50.
\end{equation}

Una volta determinato il prezzo di equilibrio, possiamo sostituirlo
in una delle due equazioni. Utilizzando, ad esempio, la funzione di
domanda otteniamo

\begin{equation}
    Q^* = 1000 - 10(50) = 500.
\end{equation}

La soluzione del modello è pertanto

\begin{equation}
    P^* = 50,
    \qquad
    Q^* = 500.
\end{equation}

Il punto $(P^*,Q^*)$ rappresenta l'equilibrio del mercato.
Graficamente, esso corrisponde al punto di intersezione tra la curva
di domanda e la curva di offerta.

Per prezzi differenti da $P^*$, domanda e offerta non coincidono.
Se il prezzo è inferiore al prezzo di equilibrio, la quantità
domandata è superiore alla quantità offerta e si verifica un
\emph{eccesso di domanda}. Se, al contrario, il prezzo è superiore
al prezzo di equilibrio, la quantità offerta è superiore alla
quantità domandata e si verifica un \emph{eccesso di offerta}.


\subsection{Dall'equilibrio al benessere sociale}

La determinazione di $P^*$ e $Q^*$ non esaurisce necessariamente
l'analisi economica. Una volta conosciuto l'equilibrio possiamo,
ad esempio, valutare il benessere generato dal mercato attraverso
il surplus dei consumatori e il surplus dei produttori.

Il \emph{surplus del consumatore} misura la differenza tra quanto
i consumatori sarebbero disposti a pagare per le unità acquistate
e quanto effettivamente pagano al prezzo di equilibrio.

Il \emph{surplus del produttore} misura invece la differenza tra
quanto i produttori ricevono dalla vendita delle unità prodotte
e il prezzo minimo al quale sarebbero disposti a offrirle.

Possiamo quindi definire il benessere sociale come

\begin{equation}
    SW = CS + PS,
    \label{eq:social_welfare}
\end{equation}

dove $CS$ indica il surplus dei consumatori e $PS$ il surplus
dei produttori.

Nel caso di funzioni lineari, questi surplus possono essere
rappresentati graficamente come aree comprese tra le curve di
domanda e offerta e il prezzo di equilibrio.

Questo passaggio mette in evidenza un aspetto importante
dell'analisi computazionale: la soluzione del modello non
rappresenta necessariamente il risultato finale dell'analisi.
Le variabili di equilibrio ottenute dalla soluzione del modello
possono diventare, a loro volta, input per il calcolo di ulteriori
grandezze di interesse economico.


\subsection{Esperimenti sul modello}

Una volta costruito e risolto il modello, possiamo utilizzarlo per
studiare come cambia l'equilibrio in seguito a una modifica dei
parametri.

Supponiamo, ad esempio, che l'intercetta della funzione di domanda
si riduca da

\begin{equation}
    \bar{Q}_d = 1000
\end{equation}

a

\begin{equation}
    \bar{Q}_d = 750.
\end{equation}

La funzione di domanda diventa

\begin{equation}
    Q_d = 750 - 10P.
\end{equation}

Possiamo quindi risolvere nuovamente il modello, determinare il
nuovo prezzo e la nuova quantità di equilibrio e, successivamente,
ricalcolare il surplus dei consumatori, il surplus dei produttori
e il benessere sociale.

Questo semplice esperimento introduce una delle idee fondamentali
dell'economia computazionale. Una volta specificato un modello,
possiamo assegnare valori ai suoi parametri, calcolarne la soluzione
e ripetere il calcolo modificando uno o più parametri per studiare
come cambiano i risultati.


\subsection{Dal modello economico al sistema di equazioni}

Nel semplice esempio precedente il prezzo e la quantità di equilibrio
possono essere determinati facilmente attraverso pochi passaggi
algebrici. Questa strategia diventa tuttavia progressivamente meno
pratica quando aumenta il numero delle equazioni e delle variabili
del modello.

È quindi utile introdurre una rappresentazione più generale del
problema.

Utilizzando la condizione di equilibrio, le equazioni
\eqref{eq:domanda} e \eqref{eq:offerta} possono essere scritte come

\begin{equation}
    Q = \bar{Q}_d - aP
\end{equation}

e

\begin{equation}
    Q = \bar{Q}_s + bP.
\end{equation}

Portando le variabili endogene a sinistra e i termini noti a destra
otteniamo

\begin{equation}
    Q + aP = \bar{Q}_d
\end{equation}

e

\begin{equation}
    Q - bP = \bar{Q}_s.
\end{equation}

Le due equazioni possono essere raccolte in un'unica espressione
matriciale:

\begin{equation}
    \begin{pmatrix}
        1 & a \\
        1 & -b
    \end{pmatrix}
    \begin{pmatrix}
        Q \\
        P
    \end{pmatrix}
    =
    \begin{pmatrix}
        \bar{Q}_d \\
        \bar{Q}_s
    \end{pmatrix}.
    \label{eq:domanda_offerta_matrice}
\end{equation}

Definiamo ora

\begin{equation}
    A =
    \begin{pmatrix}
        1 & a \\
        1 & -b
    \end{pmatrix},
    \qquad
    x =
    \begin{pmatrix}
        Q \\
        P
    \end{pmatrix},
    \qquad
    d =
    \begin{pmatrix}
        \bar{Q}_d \\
        \bar{Q}_s
    \end{pmatrix}.
\end{equation}

Il modello può quindi essere scritto nella forma compatta

\begin{equation}
    \boxed{Ax=d}.
    \label{eq:sistema_lineare_generale}
\end{equation}

La rappresentazione matriciale non modifica il contenuto economico
del modello. Il prodotto matrice-vettore $Ax$ ricostruisce infatti,
riga per riga, le equazioni originarie.

Il vantaggio della rappresentazione \eqref{eq:sistema_lineare_generale}
è che permette di riconoscere nel nostro semplice modello di domanda
e offerta un problema matematico molto più generale: la soluzione
di un sistema di equazioni lineari.


\subsection{Il problema computazionale}

Il modello di domanda e offerta ci ha quindi condotto a un problema
che incontreremo ripetutamente nel corso.

Data una matrice $A$ e un vettore $d$, vogliamo determinare il
vettore $x$ tale che

\begin{equation}
    Ax=d.
\end{equation}

Nel nostro esempio, $x$ contiene soltanto due elementi, il prezzo
e la quantità di equilibrio. Nei modelli più complessi, tuttavia,
$x$ può contenere decine, centinaia o anche migliaia di incognite.

Diventa allora necessario disporre di procedure sistematiche ed
efficienti per risolvere il sistema.

Prima di studiare queste procedure è opportuno richiamare alcuni
concetti dell'algebra matriciale che saranno utilizzati nella
costruzione degli algoritmi numerici. Essi costituiscono l'oggetto
del prossimo paragrafo.

\section{Strumenti di algebra matriciale}

Nel paragrafo precedente abbiamo visto che un semplice modello di
domanda e offerta può essere rappresentato mediante un sistema di
equazioni lineari della forma

\begin{equation}
    Ax=b.
    \label{eq:sistema_lineare}
\end{equation}

Questa rappresentazione sarà utilizzata ripetutamente nel corso.
Prima di affrontare i metodi numerici che consentono di calcolare
la soluzione di un sistema lineare, è quindi utile richiamare alcuni
concetti fondamentali dell'algebra matriciale.

L'obiettivo di questo capitolo non è fornire una trattazione completa
dell'algebra lineare, ma introdurre gli strumenti che saranno
necessari per comprendere i metodi numerici presentati nei capitoli
successivi.


\subsection{Matrici e vettori}

Una matrice è un insieme di numeri organizzati in righe e colonne.
Una matrice con $m$ righe e $n$ colonne può essere rappresentata come

\begin{equation}
A =
\begin{pmatrix}
a_{11} & a_{12} & \cdots & a_{1n} \\
a_{21} & a_{22} & \cdots & a_{2n} \\
\vdots & \vdots & \ddots & \vdots \\
a_{m1} & a_{m2} & \cdots & a_{mn}
\end{pmatrix}.
\end{equation}

L'elemento $a_{ij}$ indica l'elemento che si trova nella riga $i$
e nella colonna $j$.

Quando il numero delle righe coincide con quello delle colonne,
la matrice viene detta \emph{quadrata}. Una matrice quadrata di
dimensione $n$ appartiene all'insieme delle matrici
$\mathbb{R}^{n\times n}$.

Alcune particolari matrici quadrate avranno un ruolo importante
nei metodi numerici che studieremo.


\subsubsection{Matrice diagonale}

Una matrice è detta \emph{diagonale} quando tutti gli elementi
al di fuori della diagonale principale sono nulli:

\begin{equation}
D =
\begin{pmatrix}
d_{11} & 0      & \cdots & 0 \\
0      & d_{22} & \cdots & 0 \\
\vdots & \vdots & \ddots & \vdots \\
0      & 0      & \cdots & d_{nn}
\end{pmatrix}.
\end{equation}


\subsubsection{Matrici triangolari}

Una matrice è \emph{triangolare superiore} quando tutti gli
elementi al di sotto della diagonale principale sono nulli:

\begin{equation}
U =
\begin{pmatrix}
u_{11} & u_{12} & \cdots & u_{1n} \\
0      & u_{22} & \cdots & u_{2n} \\
\vdots & \ddots & \ddots & \vdots \\
0      & \cdots & 0      & u_{nn}
\end{pmatrix}.
\end{equation}

Analogamente, una matrice è \emph{triangolare inferiore} quando
sono nulli tutti gli elementi al di sopra della diagonale:

\begin{equation}
L =
\begin{pmatrix}
l_{11} & 0      & \cdots & 0 \\
l_{21} & l_{22} & \cdots & 0 \\
\vdots & \vdots & \ddots & \vdots \\
l_{n1} & l_{n2} & \cdots & l_{nn}
\end{pmatrix}.
\end{equation}

Queste matrici avranno particolare importanza quando introdurremo
i metodi numerici per la soluzione dei sistemi lineari.


\subsubsection{Matrice identità}

Un caso particolare di matrice diagonale è la matrice identità,

\begin{equation}
I =
\begin{pmatrix}
1 & 0 & \cdots & 0 \\
0 & 1 & \cdots & 0 \\
\vdots & \vdots & \ddots & \vdots \\
0 & 0 & \cdots & 1
\end{pmatrix}.
\end{equation}

La matrice identità svolge nel prodotto tra matrici un ruolo
analogo a quello svolto dal numero $1$ nella moltiplicazione
tra numeri. Per una matrice $A$ di dimensioni compatibili vale

\begin{equation}
AI=IA=A.
\end{equation}


\subsection{Matrice trasposta e matrici simmetriche}

La \emph{trasposta} di una matrice $A$ si ottiene scambiando
le righe con le colonne. Essa viene indicata con $A^T$.

Se, ad esempio,

\begin{equation}
A =
\begin{pmatrix}
a & b \\
c & d
\end{pmatrix},
\end{equation}

allora

\begin{equation}
A^T =
\begin{pmatrix}
a & c \\
b & d
\end{pmatrix}.
\end{equation}

La trasposta di una matrice triangolare superiore è quindi una
matrice triangolare inferiore, e viceversa.

Una matrice quadrata è detta \emph{simmetrica} quando coincide
con la propria trasposta:

\begin{equation}
A=A^T.
\end{equation}

Equivalentemente, in una matrice simmetrica vale

\begin{equation}
a_{ij}=a_{ji}.
\end{equation}

La simmetria sarà una proprietà importante perché alcune tecniche
numeriche possono essere utilizzate in modo particolarmente
efficiente quando la matrice del sistema possiede questa
caratteristica.


\subsection{Operazioni tra matrici}

Per utilizzare la rappresentazione matriciale dei modelli è
necessario conoscere alcune operazioni fondamentali.


\subsubsection{Somma}

Due matrici $A$ e $B$ con le stesse dimensioni possono essere
sommate elemento per elemento. Se

\begin{equation}
C=A+B,
\end{equation}

allora

\begin{equation}
c_{ij}=a_{ij}+b_{ij}.
\end{equation}


\subsubsection{Moltiplicazione per uno scalare}

Se $d$ è un numero reale e $A$ una matrice, il prodotto

\begin{equation}
C=dA
\end{equation}

si ottiene moltiplicando ogni elemento di $A$ per $d$:

\begin{equation}
c_{ij}=d\,a_{ij}.
\end{equation}


\subsubsection{Prodotto tra matrici}

Il prodotto tra matrici richiede maggiore attenzione. Date due
matrici $A$ e $B$, il prodotto

\begin{equation}
C=AB
\end{equation}

è definito se il numero di colonne di $A$ coincide con il numero
di righe di $B$.

Ogni elemento della matrice risultante viene ottenuto moltiplicando
gli elementi della corrispondente riga di $A$ per quelli della
corrispondente colonna di $B$ e sommando i prodotti:

\begin{equation}
c_{ij}=\sum_z a_{iz}b_{zj}.
\label{eq:prodotto_matrici}
\end{equation}

Se $A$ ha dimensione $m\times n$ e $B$ ha dimensione $n\times p$,
la matrice risultante $C=AB$ avrà dimensione $m\times p$.

È importante osservare che, in generale,

\begin{equation}
AB \neq BA.
\end{equation}

Il prodotto tra matrici non gode quindi della proprietà commutativa.

Questa regola di moltiplicazione spiega perché il sistema

\begin{equation}
Ax=b
\end{equation}

rappresenta simultaneamente tutte le equazioni di un modello.
Ogni riga della matrice $A$, moltiplicata per il vettore $x$,
genera infatti una delle equazioni del sistema.


\subsection{Alcune grandezze associate a una matrice}

A una matrice quadrata possono essere associate alcune grandezze
scalari. Tra queste consideriamo la traccia e il determinante.


\subsubsection{Traccia}

La \emph{traccia} di una matrice quadrata è la somma degli elementi
della diagonale principale:

\begin{equation}
\operatorname{tr}(A)
=
\sum_{i=1}^{n}a_{ii}.
\end{equation}

Ad esempio, se

\begin{equation}
A=
\begin{pmatrix}
2 & 3 \\
4 & 5
\end{pmatrix},
\end{equation}

allora

\begin{equation}
\operatorname{tr}(A)=2+5=7.
\end{equation}


\subsubsection{Determinante}

Il \emph{determinante} associa a una matrice quadrata $A$ un
numero reale, indicato con

\begin{equation}
\det(A).
\end{equation}

Il determinante avrà un ruolo importante nello studio dei sistemi
lineari perché consente, tra le altre cose, di stabilire se una
matrice quadrata è invertibile.


\subsection{Calcolo del determinante}

Per una matrice generale di dimensione $n\times n$, il determinante
può essere calcolato mediante lo sviluppo di Laplace.

Scegliendo una qualsiasi riga $i$,

\begin{equation}
\det(A)
=
\sum_{j=1}^{n}
a_{ij}(-1)^{i+j}\det(A_{ij}),
\label{eq:laplace_riga}
\end{equation}

dove $A_{ij}$ è la matrice ottenuta eliminando da $A$ la riga $i$
e la colonna $j$.

In modo equivalente, scegliendo una colonna $j$,

\begin{equation}
\det(A)
=
\sum_{i=1}^{n}
a_{ij}(-1)^{i+j}\det(A_{ij}).
\label{eq:laplace_colonna}
\end{equation}

Il risultato non dipende dalla riga o dalla colonna scelta per
effettuare lo sviluppo.

La natura ricorsiva di questa procedura rende tuttavia il calcolo
del determinante progressivamente più laborioso all'aumentare
della dimensione della matrice. Questo aspetto sarà importante
quando passeremo dalla soluzione algebrica di piccoli sistemi
alla loro soluzione numerica.


\subsubsection{Il caso delle matrici triangolari}

Il calcolo del determinante è particolarmente semplice nel caso
delle matrici diagonali e triangolari.

Se $A$ è triangolare, vale infatti

\begin{equation}
\det(A)=\prod_{i=1}^{n}a_{ii}.
\label{eq:det_triangolare}
\end{equation}

Il determinante è quindi semplicemente il prodotto degli elementi
della diagonale principale.



\subsection{Matrici singolari e non singolari}

Una matrice quadrata $A$ è detta \emph{non singolare} quando

\begin{equation}
\det(A)\neq 0.
\end{equation}

Se invece

\begin{equation}
\det(A)=0,
\end{equation}

la matrice è detta \emph{singolare}.

Questa distinzione è fondamentale per il problema che stiamo
studiando. Una matrice $A$ possiede infatti una matrice inversa
se e solo se è non singolare.

Ricordando quanto detto al paragrafo precedente, possiamo dire che una matrice triangolare è non singolare se tutti gli elementi della sua diagonale sono diversi da zero.

\subsection{La matrice inversa}

L'inversa di una matrice quadrata $A$, indicata con $A^{-1}$,
è definita come la matrice che soddisfa

\begin{equation}
A^{-1}A=I.
\end{equation}

dove $I$ è la matrice identità definita in precedenza.

Vale inoltre

\begin{equation}
AA^{-1}=I.
\end{equation}

L'analogia con i numeri reali è immediata. Se $a\neq0$, il suo
inverso $a^{-1}=1/a$ soddisfa

\begin{equation}
a^{-1}a=1.
\end{equation}

Come ricordato in precedenza, nel caso delle matrici, il ruolo del numero $1$ è svolto dalla
matrice identità $I$.


\subsection{Soluzione di un sistema mediante la matrice inversa}

Torniamo ora al problema dal quale siamo partiti:

\begin{equation}
Ax=b.
\end{equation}

Se $A$ è non singolare, esiste $A^{-1}$. Possiamo quindi
premoltiplicare entrambi i membri dell'equazione per $A^{-1}$:

\begin{equation}
A^{-1}Ax=A^{-1}b.
\end{equation}

Poiché

\begin{equation}
A^{-1}A=I,
\end{equation}

otteniamo

\begin{equation}
Ix=A^{-1}b.
\end{equation}

Infine, dato che $Ix=x$,

\begin{equation}
\boxed{x=A^{-1}b}.
\label{eq:soluzione_inversa}
\end{equation}

Abbiamo quindi ottenuto una formula che, almeno dal punto di vista
algebrico, permette di risolvere il sistema lineare.


\subsection{Dal risultato al problema computazionale}

La relazione

\begin{equation}
x=A^{-1}b
\end{equation}

fornisce una soluzione elegante del problema dal punto di vista
matematico. Potrebbe quindi sembrare che il problema posto alla
fine del capitolo precedente sia già risolto: basta calcolare
l'inversa di $A$ e moltiplicarla per $b$.

Dal punto di vista computazionale, tuttavia, rimane una questione
fondamentale:

Il calcolo della matrice inversa è agevole dal punto di vista computazionale?

Purtroppo, il numero di operazioni che il calcolatore deve effettuare, cresce molto rapidamente all'aumentare della dimensione del sistema.

Dunque, ancorché calcolare l'inversa sia possibile per piccoli sistemi, occorre identificare strategie alternative. 

%Per sistemi di piccole dimensioni la distinzione può sembrare poco importante. Quando il numero delle equazioni cresce, invece, il numero di operazioni richieste e il modo in cui esse vengono organizzate diventano elementi centrali.

La formula $x=A^{-1}b$ deve quindi essere interpretata innanzitutto
come una caratterizzazione matematica della soluzione. Non implica
necessariamente che il calcolo esplicito dell'inversa sia il modo
più conveniente per ottenere numericamente $x$.

Questo ci conduce direttamente al problema affrontato nel prossimo
capitolo: individuare procedure più adatte per risolvere un sistema
lineare senza dover necessariamente calcolare esplicitamente
$A^{-1}$.

\section{Soluzione di sistemi lineari mediante eliminazione gaussiana}
\label{sec:eliminazione-gaussiana}

Nella sezione precedente abbiamo visto che, se la matrice $A$ è
non singolare, il sistema lineare

\begin{equation}
    Ax=b
\end{equation}

ammette una soluzione unica, che può essere espressa formalmente come

\begin{equation}
    x=A^{-1}b.
\end{equation}

Questa espressione caratterizza la soluzione dal punto di vista
matematico, ma non implica che il calcolo esplicito della matrice
inversa rappresenti il modo più conveniente per determinarla.

Un'altra possibilità, nota dall'algebra lineare, è rappresentata
dalla regola di Cramer. Per un sistema di $n$ equazioni in $n$
incognite, ciascun elemento della soluzione può essere scritto come

\begin{equation}
    x_i =
    \frac{\det(A_i)}{\det(A)},
    \qquad i=1,\ldots,n,
    \label{eq:cramer}
\end{equation}

dove $A_i$ è la matrice ottenuta sostituendo la $i$-esima colonna
di $A$ con il vettore $b$.

Anche questa formula fornisce una caratterizzazione esatta della
soluzione. Tuttavia, il suo utilizzo richiede il calcolo di numerosi
determinanti. Il costo computazionale cresce quindi molto rapidamente
all'aumentare della dimensione del sistema. La regola di Cramer è
utile per comprendere alcune proprietà teoriche dei sistemi lineari,
ma non costituisce un metodo efficiente per la soluzione numerica di
sistemi di grandi dimensioni.

Occorre quindi cercare una procedura che permetta di ottenere la
soluzione senza calcolare esplicitamente l'inversa di $A$ e senza
dover calcolare un determinante per ogni incognita.

Una delle procedure fondamentali per raggiungere questo obiettivo è
l'\emph{eliminazione gaussiana}.


\subsection{L'idea dell'eliminazione gaussiana}

L'idea alla base dell'eliminazione gaussiana consiste nel trasformare
il sistema originario in un sistema equivalente, ma più semplice da
risolvere.

In particolare, si cerca di trasformare la matrice del sistema in una
matrice triangolare. Come vedremo, un sistema caratterizzato da una
matrice triangolare può essere risolto in modo particolarmente
semplice, procedendo sequenzialmente da un'equazione all'altra.

Per descrivere contemporaneamente i coefficienti del sistema e i
termini noti, è conveniente costruire la \emph{matrice aumentata}

\begin{equation}
    C = [A \mid b].
    \label{eq:matrice-aumentata}
\end{equation}

L'eliminazione gaussiana modifica progressivamente questa matrice
attraverso operazioni che non modificano la soluzione del sistema.
In questo modo otteniamo una successione di sistemi equivalenti,
fino a raggiungere una forma che può essere risolta facilmente.


\subsection{Operazioni elementari sulle righe}

Le trasformazioni utilizzate nell'eliminazione gaussiana sono
costituite da tre operazioni elementari sulle righe della matrice
aumentata:

\begin{enumerate}
    \item scambiare la posizione di due righe;
    
    \item moltiplicare tutti gli elementi di una riga per uno
    scalare non nullo;
    
    \item sostituire una riga con la somma della stessa riga e
    di un multiplo di un'altra riga.
\end{enumerate}

Queste operazioni trasformano il sistema in un sistema equivalente:
cambiano cioè la forma nella quale le equazioni sono rappresentate,
ma non l'insieme delle loro soluzioni.

La terza operazione è particolarmente importante. Scegliendo
opportunamente il coefficiente con il quale moltiplicare una riga,
possiamo infatti annullare uno degli elementi della matrice.
Ripetendo questa operazione possiamo progressivamente ottenere una
matrice triangolare.


\subsection{Applicazione al modello di domanda e offerta}

Riprendiamo il modello di domanda e offerta introdotto nella
Sezione~\ref{sec:domanda-offerta}. Il sistema è

\begin{equation}
    Ax=b,
\end{equation}

con

\begin{equation}
    A=
    \begin{pmatrix}
        1 & a\\
        1 & -b
    \end{pmatrix},
    \qquad
    x=
    \begin{pmatrix}
        Q\\
        P
    \end{pmatrix},
    \qquad
    b=
    \begin{pmatrix}
        \bar Q_d\\
        \bar Q_s
    \end{pmatrix}.
\end{equation}

Utilizzando i valori

\begin{equation}
    \bar Q_d=1000,
    \qquad
    \bar Q_s=250,
    \qquad
    a=10,
    \qquad
    b=5,
\end{equation}

otteniamo

\begin{equation}
    \begin{pmatrix}
        1 & 10\\
        1 & -5
    \end{pmatrix}
    \begin{pmatrix}
        Q\\
        P
    \end{pmatrix}
    =
    \begin{pmatrix}
        1000\\
        250
    \end{pmatrix}.
    \label{eq:sistema-domanda-offerta-gauss}
\end{equation}

La corrispondente matrice aumentata è

\begin{equation}
    C=[A\mid b]
    =
    \left(
    \begin{array}{cc|c}
        1 & 10 & 1000\\
        1 & -5 & 250
    \end{array}
    \right).
    \label{eq:matrice-aumentata-domanda-offerta}
\end{equation}

Il nostro obiettivo è trasformare questa matrice in modo da
eliminare una delle due incognite da una delle equazioni.


\subsection{Una prima trasformazione}

Una prima possibilità consiste nel sostituire la prima riga con la
somma della prima riga e due volte la seconda riga. Indicando le
righe con $R_1$ e $R_2$, l'operazione può essere scritta come

\begin{equation}
    R_1 \leftarrow R_1+2R_2.
\end{equation}

Applicando l'operazione alla matrice aumentata otteniamo

\begin{equation}
    \left(
    \begin{array}{cc|c}
        1 & 10 & 1000\\
        1 & -5 & 250
    \end{array}
    \right)
    \longrightarrow
    \left(
    \begin{array}{cc|c}
        3 & 0 & 1500\\
        1 & -5 & 250
    \end{array}
    \right).
\end{equation}

La prima equazione del nuovo sistema è quindi

\begin{equation}
    3Q=1500,
\end{equation}

da cui otteniamo immediatamente

\begin{equation}
    Q=500.
\end{equation}

Conosciuto $Q$, possiamo utilizzare la seconda equazione,

\begin{equation}
    Q-5P=250,
\end{equation}

e sostituire il valore appena calcolato:

\begin{equation}
    500-5P=250.
\end{equation}

Da cui

\begin{equation}
    P=50.
\end{equation}

Otteniamo pertanto nuovamente la soluzione

\begin{equation}
    x^*=
    \begin{pmatrix}
        500\\
        50
    \end{pmatrix}.
\end{equation}


\subsection{Una seconda trasformazione}

La sequenza di operazioni utilizzata per risolvere un sistema non è
necessariamente unica. Nel nostro semplice esempio possiamo, ad
esempio, scegliere di sostituire la prima riga con la differenza tra
la prima e la seconda riga:

\begin{equation}
    R_1 \leftarrow R_1-R_2.
\end{equation}

Otteniamo

\begin{equation}
    \left(
    \begin{array}{cc|c}
        1 & 10 & 1000\\
        1 & -5 & 250
    \end{array}
    \right)
    \longrightarrow
    \left(
    \begin{array}{cc|c}
        0 & 15 & 750\\
        1 & -5 & 250
    \end{array}
    \right).
\end{equation}

La prima equazione è ora

\begin{equation}
    15P=750,
\end{equation}

da cui

\begin{equation}
    P=50.
\end{equation}

Sostituendo questo valore nella seconda equazione,

\begin{equation}
    Q-5P=250,
\end{equation}

otteniamo

\begin{equation}
    Q-5(50)=250,
\end{equation}

e quindi

\begin{equation}
    Q=500.
\end{equation}

Anche questa sequenza di trasformazioni conduce dunque alla stessa
soluzione:

\begin{equation}
    Q^*=500,
    \qquad
    P^*=50.
\end{equation}


\subsection{Dall'esempio all'algoritmo}

Il sistema di domanda e offerta contiene soltanto due equazioni e
una singola operazione di eliminazione è sufficiente per renderlo
immediatamente risolvibile. Tuttavia, la stessa idea può essere
applicata a sistemi con un numero maggiore di equazioni.

In termini generali, l'eliminazione gaussiana utilizza le operazioni
elementari sulle righe per eliminare progressivamente le incognite
dalle diverse equazioni. L'obiettivo è trasformare il sistema
originario

\begin{equation}
    Ax=b
\end{equation}

in un sistema equivalente

\begin{equation}
    Ux=c,
\end{equation}

nel quale $U$ è una matrice triangolare superiore.

La trasformazione non modifica la soluzione del problema, ma produce
un sistema la cui struttura rende molto più semplice il calcolo delle
incognite.

Una volta raggiunta la forma triangolare, la soluzione può essere
determinata partendo dall'ultima equazione e procedendo verso la
prima. Questa procedura prende il nome di \emph{sostituzione
all'indietro}.

Nel caso di una matrice triangolare inferiore si può procedere in
modo analogo, partendo dalla prima equazione. La procedura prende
in questo caso il nome di \emph{sostituzione in avanti}.


\subsection{Soluzione mediante software}

Nelle applicazioni non è normalmente necessario eseguire
manualmente tutte le operazioni di eliminazione. I principali
software per il calcolo scientifico dispongono di funzioni
specifiche per la soluzione di sistemi lineari.

Ad esempio, in \textsf{R} il sistema può essere risolto mediante

\begin{verbatim}
A <- matrix(c(1,10,1,-5), ncol=2, byrow=TRUE)
b <- c(1000,250)
solve(A,b)
\end{verbatim}

In \textsf{Python}, utilizzando \textsf{NumPy}, possiamo scrivere

\begin{verbatim}
import numpy as np

A = [[1,10],
     [1,-5]]

b = [[1000],
     [250]]

np.linalg.solve(A,b)
\end{verbatim}

mentre in \textsf{Julia} e \textsf{Octave} il sistema può essere
risolto attraverso l'operatore \texttt{\textbackslash}:

\begin{verbatim}
A = [1 10; 1 -5]
b = [1000; 250]
A\b
\end{verbatim}

In tutti i casi il risultato è

\begin{equation}
    x^*=
    \begin{pmatrix}
        500\\
        50
    \end{pmatrix}.
\end{equation}

L'utilizzo di una funzione software nasconde all'utente molte delle
operazioni numeriche necessarie per ottenere la soluzione. Comprendere
l'eliminazione gaussiana rimane tuttavia importante, perché permette
di capire la logica sulla quale si basano i metodi più efficienti
che utilizzeremo successivamente.

\section{Il modello di domanda e offerta con scambi tra due paesi}
\label{sec:commercio-due-paesi}

Il modello di domanda e offerta introdotto nella
Sezione~\ref{sec:domanda-offerta} può essere esteso considerando
due paesi che producono e consumano lo stesso bene e che hanno
la possibilità di commerciare tra loro.

Questa estensione permette di utilizzare gli strumenti introdotti
nelle sezioni precedenti per studiare un problema economico più
articolato. In particolare, vogliamo determinare la direzione degli
scambi, il prezzo che si forma in presenza di commercio
internazionale e la quantità scambiata tra i due paesi.

Successivamente utilizzeremo il modello per studiare gli effetti
di due interventi di politica commerciale: una tariffa sulle
importazioni e un sussidio alle esportazioni.


\subsection{I mercati nazionali}

Consideriamo due paesi, che indichiamo con $H$ (\emph{Home}) e
$F$ (\emph{Foreign}).

Nel paese $H$, domanda e offerta sono descritte dalle funzioni

\begin{equation}
    Q_{d,H} = \bar{Q}_{d,H} - a_H P_H,
    \label{eq:domanda-H}
\end{equation}

\begin{equation}
    Q_{s,H} = \bar{Q}_{s,H} + b_H P_H.
    \label{eq:offerta-H}
\end{equation}

Analogamente, nel paese $F$ abbiamo

\begin{equation}
    Q_{d,F} = \bar{Q}_{d,F} - a_F P_F,
    \label{eq:domanda-F}
\end{equation}

\begin{equation}
    Q_{s,F} = \bar{Q}_{s,F} + b_F P_F.
    \label{eq:offerta-F}
\end{equation}

Se i due paesi non possono commerciare tra loro, ciascun mercato
raggiunge autonomamente il proprio equilibrio. Possiamo quindi
calcolare un prezzo di equilibrio per $H$, che indichiamo con
$P_{e,H}$, e un prezzo di equilibrio per $F$, indicato con
$P_{e,F}$.

Il confronto tra questi due prezzi permette di determinare quale
paese avrà convenienza a importare il bene e quale avrà convenienza
a esportarlo.


\subsection{Dall'autarchia al commercio internazionale}

Supponiamo che il prezzo di equilibrio in $H$ sia maggiore del
prezzo di equilibrio in $F$:

\begin{equation}
    P_{e,H} > P_{e,F}.
    \label{eq:prezzi-autarchia}
\end{equation}

Prima dell'apertura al commercio, il bene è quindi relativamente
più costoso in $H$ e relativamente meno costoso in $F$.

Se diventa possibile commerciare tra i due paesi, i consumatori
di $H$ hanno convenienza ad acquistare il bene prodotto in $F$.
Allo stesso tempo, i produttori di $F$ hanno convenienza a vendere
parte della propria produzione nel mercato di $H$.

Il paese $H$ diventa quindi importatore, mentre il paese $F$
diventa esportatore.

Per determinare l'equilibrio internazionale dobbiamo costruire
una domanda e un'offerta sul mercato mondiale a partire dalle
curve di domanda e offerta dei due mercati nazionali.


\subsection{La domanda di importazioni}

La quantità che il paese $H$ desidera importare corrisponde alla
differenza tra la quantità domandata e quella prodotta internamente.
Definiamo quindi la domanda di importazioni di $H$ come

\begin{equation}
    M_H = Q_{d,H} - Q_{s,H}.
\end{equation}

Sostituendo le funzioni di domanda e offerta otteniamo

\begin{align}
    M_H
    &=
    \bar{Q}_{d,H}-a_HP_H
    -
    \left(\bar{Q}_{s,H}+b_HP_H\right)
    \\
    &=
    \left(\bar{Q}_{d,H}-\bar{Q}_{s,H}\right)
    -
    \left(a_H+b_H\right)P_H.
\end{align}

Possiamo semplificare la notazione definendo

\begin{equation}
    \bar{Q}_{d,W}
    =
    \bar{Q}_{d,H}-\bar{Q}_{s,H}
\end{equation}

e

\begin{equation}
    a_W=a_H+b_H.
\end{equation}

La domanda sul mercato internazionale può quindi essere scritta
come

\begin{equation}
    M_H=\bar{Q}_{d,W}-a_WP_W.
    \label{eq:domanda-mondiale}
\end{equation}

Essa è decrescente rispetto al prezzo: all'aumentare del prezzo
internazionale, la differenza tra domanda e offerta interna di $H$
si riduce e, di conseguenza, diminuisce la quantità che il paese
desidera importare.


\subsection{L'offerta di esportazioni}

Il paese $F$, al contrario, può esportare la parte della propria
produzione che eccede la domanda interna.

Definiamo quindi l'offerta di esportazioni come

\begin{equation}
    E_F=Q_{s,F}-Q_{d,F}.
\end{equation}

Sostituendo le rispettive funzioni otteniamo

\begin{align}
    E_F
    &=
    \bar{Q}_{s,F}+b_FP_F
    -
    \left(\bar{Q}_{d,F}-a_FP_F\right)
    \\
    &=
    \left(\bar{Q}_{s,F}-\bar{Q}_{d,F}\right)
    +
    \left(b_F+a_F\right)P_F.
\end{align}

Definiamo

\begin{equation}
    \bar{Q}_{s,W}
    =
    \bar{Q}_{s,F}-\bar{Q}_{d,F}
\end{equation}

e

\begin{equation}
    b_W=b_F+a_F.
\end{equation}

L'offerta sul mercato internazionale diventa quindi

\begin{equation}
    E_F=\bar{Q}_{s,W}+b_WP_W.
    \label{eq:offerta-mondiale}
\end{equation}

A differenza della domanda di importazioni, l'offerta di
esportazioni è crescente rispetto al prezzo.


\subsection{L'equilibrio sul mercato internazionale}

In equilibrio, la quantità che $H$ desidera importare deve
coincidere con la quantità che $F$ desidera esportare:

\begin{equation}
    M_H=E_F=Q_W.
\end{equation}

Il problema può pertanto essere scritto come

\begin{equation}
    Q_W=\bar{Q}_{d,W}-a_WP_W,
\end{equation}

\begin{equation}
    Q_W=\bar{Q}_{s,W}+b_WP_W.
\end{equation}

Ritroviamo quindi esattamente la stessa struttura del modello di
domanda e offerta studiato nella Sezione~\ref{sec:domanda-offerta}.
Anche il mercato internazionale può essere rappresentato mediante
un sistema lineare:

\begin{equation}
    \begin{pmatrix}
        1 & a_W\\
        1 & -b_W
    \end{pmatrix}
    \begin{pmatrix}
        Q_W\\
        P_W
    \end{pmatrix}
    =
    \begin{pmatrix}
        \bar{Q}_{d,W}\\
        \bar{Q}_{s,W}
    \end{pmatrix}.
    \label{eq:sistema-mercato-internazionale}
\end{equation}

La soluzione del sistema determina la quantità scambiata
$Q_W^*$ e il prezzo internazionale $P_W^*$.

In assenza di costi o impedimenti al commercio, il prezzo
internazionale vale per entrambi i paesi:

\begin{equation}
    P_H=P_F=P_W^*.
\end{equation}

Al prezzo $P_W^*$, nel paese $H$ la domanda supera la produzione
nazionale esattamente di $Q_W^*$, mentre nel paese $F$ la produzione
supera la domanda interna della stessa quantità. Il paese $F$
esporta pertanto $Q_W^*$ unità del bene verso il paese $H$.


\subsection{Introduzione di una tariffa}

Supponiamo ora che il paese $H$ voglia limitare le importazioni
introducendo una tariffa $t$ per ogni unità del bene importata da
$F$.

In assenza della tariffa, un produttore di $F$ può vendere il bene
nel proprio paese oppure in $H$ allo stesso prezzo internazionale.
Con l'introduzione della tariffa, invece, l'esportatore deve
sostenere un costo aggiuntivo pari a $t$ per ogni unità venduta
in $H$.

Immediatamente dopo l'introduzione della tariffa, esportare verso
$H$ diventa quindi meno conveniente. La riduzione degli scambi
produce due effetti.

Nel paese $H$ si genera un eccesso di domanda e il prezzo tende
ad aumentare. Nel paese $F$ si genera invece un eccesso di offerta
e il prezzo tende a diminuire.

Il processo continua fino a quando la differenza tra i prezzi nei
due paesi è esattamente pari alla tariffa:

\begin{equation}
    P_H=P_F+t.
    \label{eq:tariffa-prezzi}
\end{equation}

La tariffa crea quindi un divario tra i prezzi nazionali e riduce
la quantità scambiata rispetto all'equilibrio di libero commercio.


\subsection{Introduzione di un sussidio alle esportazioni}

Consideriamo ora una politica differente. Supponiamo che il paese
$F$ eroghi ai propri produttori un sussidio $t$ per ogni unità
esportata verso $H$.

Il sussidio aumenta la convenienza a vendere nel mercato estero.
Le imprese di $F$ cercano quindi di aumentare le esportazioni verso
$H$.

Nel paese $F$ si genera una pressione al rialzo sul prezzo interno,
mentre in $H$ la maggiore disponibilità di beni importati esercita
una pressione al ribasso sul prezzo.

Il processo continua fino a quando

\begin{equation}
    P_F=P_H+t.
    \label{eq:sussidio-prezzi}
\end{equation}

A differenza della tariffa, il sussidio determina quindi una
quantità scambiata superiore a quella che si avrebbe in assenza
dell'intervento.


\subsection{Dal modello economico a una procedura di calcolo}

Possiamo ora tradurre il ragionamento economico in una sequenza
ordinata di operazioni.

Date le curve di domanda e offerta dei due paesi, la procedura
consiste innanzitutto nel costruire le curve di domanda e offerta
sul mercato internazionale e nel determinarne l'equilibrio.

Indichiamo con $Q_{\mathrm{ex}}$ la quantità che viene effettivamente
scambiata dopo l'intervento. Nel caso di una tariffa scegliamo una
quantità inferiore alla quantità di libero scambio:

\begin{equation}
    Q_{\mathrm{ex}}<Q_W^*,
\end{equation}

mentre nel caso di un sussidio consideriamo

\begin{equation}
    Q_{\mathrm{ex}}>Q_W^*.
\end{equation}

Una volta determinata $Q_{\mathrm{ex}}$, il prezzo nel paese $H$
può essere ricavato dalla curva di domanda internazionale, mentre
il prezzo nel paese $F$ può essere ricavato dalla curva di offerta
internazionale.

La differenza tra i due prezzi determina l'ammontare della tariffa
o del sussidio:

\begin{equation}
    \mathrm{priceGap}
    =
    \left|P_H-P_F\right|.
    \label{eq:price-gap}
\end{equation}

Abbiamo così trasformato il modello teorico in una vera e propria
\emph{procedura di calcolo}: a partire da determinati dati di input,
eseguiamo una sequenza di operazioni per ottenere gli output di
interesse.


\subsection{Valutazione degli effetti sul benessere}

Conoscere i nuovi prezzi e le quantità scambiate permette di
studiare gli effetti delle politiche commerciali sul benessere
dei due paesi.

Come nel modello di un singolo mercato, consideriamo il surplus
dei consumatori e quello dei produttori. In questo caso dobbiamo
però calcolarli separatamente per $H$ e per $F$, utilizzando i
rispettivi prezzi nazionali.


\subsubsection{Surplus dei consumatori}

Consideriamo genericamente uno dei due paesi e indichiamo con
$P_{\mathrm{naz}}$ il prezzo nazionale.

Utilizzando la curva di domanda

\begin{equation}
    Q_d=\bar Q_d-aP,
\end{equation}

determiniamo innanzitutto la quantità domandata al prezzo
$P_{\mathrm{naz}}$:

\begin{equation}
    Q_{d,\mathrm{naz}}
    =
    \bar Q_d-aP_{\mathrm{naz}}.
\end{equation}

Indicando con $P_{\max}$ l'intercetta della curva di domanda
sull'asse dei prezzi, il surplus dei consumatori è

\begin{equation}
    CS
    =
    \frac{
        Q_{d,\mathrm{naz}}
        \left(P_{\max}-P_{\mathrm{naz}}\right)
    }{2}.
    \label{eq:consumer-surplus-internazionale}
\end{equation}


\subsubsection{Surplus dei produttori}

Utilizzando la curva di offerta

\begin{equation}
    Q_s=\bar Q_s+bP,
\end{equation}

calcoliamo la quantità offerta al prezzo nazionale:

\begin{equation}
    Q_{s,\mathrm{naz}}
    =
    \bar Q_s+bP_{\mathrm{naz}}.
\end{equation}

Quando $\bar Q_s>0$, il surplus dei produttori può essere
calcolato come

\begin{equation}
    PS
    =
    \bar Q_s P_{\mathrm{naz}}
    +
    \frac{
        \left(Q_{s,\mathrm{naz}}-\bar Q_s\right)
        P_{\mathrm{naz}}
    }{2}.
    \label{eq:producer-surplus-positive-intercept}
\end{equation}

Se invece l'intercetta quantitativa dell'offerta non è positiva,
possiamo individuare l'intercetta della curva di offerta sull'asse
dei prezzi, $P_{\min}$, e calcolare

\begin{equation}
    PS
    =
    \frac{
        Q_{s,\mathrm{naz}}
        \left(P_{\mathrm{naz}}-P_{\min}\right)
    }{2}.
    \label{eq:producer-surplus-negative-intercept}
\end{equation}


\subsection{Il surplus del governo}

L'introduzione di una tariffa o di un sussidio coinvolge anche il
governo.

Nel caso della tariffa, il governo di $H$ riceve un gettito pari
alla tariffa unitaria moltiplicata per la quantità importata:

\begin{equation}
    GS_H
    =
    \mathrm{priceGap}\,Q_{\mathrm{ex}},
\end{equation}

mentre

\begin{equation}
    GS_F=0.
\end{equation}

Nel caso del sussidio, il governo di $F$ sostiene invece una spesa.
Possiamo rappresentarla come un surplus negativo:

\begin{equation}
    GS_F
    =
    -\mathrm{priceGap}\,Q_{\mathrm{ex}},
\end{equation}

mentre

\begin{equation}
    GS_H=0.
\end{equation}

Per ciascun paese il surplus totale è quindi dato dalla somma di
surplus dei consumatori, surplus dei produttori e surplus del
governo:

\begin{equation}
    TS=CS+PS+GS.
    \label{eq:total-surplus-trade}
\end{equation}


\subsection{Confrontare scenari alternativi}

Per valutare l'effetto di una politica commerciale non è sufficiente
calcolare il surplus dopo l'intervento. È necessario confrontarlo
con una situazione di riferimento.

Nel caso di una tariffa, ad esempio, dobbiamo calcolare per ciascun
paese il surplus totale prima e dopo la sua introduzione.

Per il paese $H$ avremo

\begin{equation}
    \Delta TS_H
    =
    TS_H^{\mathrm{dopo}}
    -
    TS_H^{\mathrm{prima}},
\end{equation}

mentre per il paese $F$

\begin{equation}
    \Delta TS_F
    =
    TS_F^{\mathrm{dopo}}
    -
    TS_F^{\mathrm{prima}}.
\end{equation}

Il confronto tra questi valori permette di valutare come
l'intervento modifica il benessere nei due paesi.

Questo esempio evidenzia una caratteristica tipica dell'analisi
computazionale. Una stessa procedura, come il calcolo dei surplus,
deve essere eseguita ripetutamente: per paesi diversi, per prezzi
diversi e per scenari diversi.

Riscrivere ogni volta le stesse operazioni sarebbe inefficiente
e aumenterebbe la possibilità di commettere errori. È quindi utile
organizzare la procedura di calcolo in modo che possa essere
riutilizzata.


\subsection{Funzioni definite dall'utente}

Una \emph{funzione} permette di raccogliere una sequenza di
operazioni che deve essere eseguita più volte.

Una funzione riceve uno o più valori di \emph{input}, li elabora
attraverso una procedura definita dall'utente e restituisce uno o
più valori di \emph{output}.

Nel nostro caso, una procedura generale per il calcolo dei surplus
può ricevere come input:

\begin{itemize}
    \item il prezzo del bene;
    \item i parametri della curva di domanda;
    \item i parametri della curva di offerta.
\end{itemize}

Utilizzando questi valori, la funzione calcola le quantità domandate
e offerte e restituisce come output il surplus dei consumatori e il
surplus dei produttori.

In forma schematica possiamo quindi rappresentare la funzione come

\begin{equation}
\boxed{
\text{input}
\longrightarrow
\text{procedura di calcolo}
\longrightarrow
\text{output}
}
\end{equation}

In \textsf{R}, la struttura generale di una funzione definita
dall'utente è

\begin{verbatim}
nomeFunzione <- function(input1, input2, ...) {

    # operazioni sugli input

    return(output)
}
\end{verbatim}

Una volta definita, la stessa funzione può essere richiamata
assegnando valori differenti agli input:

\begin{verbatim}
output <- nomeFunzione(valore1, valore2, ...)
\end{verbatim}

Il vantaggio non consiste soltanto nel ridurre la quantità di
codice. La funzione separa una procedura generale dai valori
specifici sui quali viene applicata. Possiamo quindi utilizzare
lo stesso algoritmo per analizzare paesi, prezzi e scenari
differenti senza dover riscrivere ogni volta la procedura di
calcolo.


\subsection{Dal problema economico all'algoritmo}

L'estensione del modello di domanda e offerta a due paesi permette
di mettere in evidenza una distinzione importante.

Il \emph{modello economico} specifica le relazioni tra domanda,
offerta, prezzi e quantità e descrive le condizioni che
caratterizzano l'equilibrio. L'\emph{algoritmo} specifica invece
la sequenza di operazioni attraverso la quale, a partire dai
parametri del modello, il computer determina i risultati.

Nel problema considerato in questa sezione possiamo sintetizzare
il procedimento come

\begin{equation}
\boxed{
\begin{array}{c}
\text{domanda e offerta nazionali}\\
\downarrow\\
\text{prezzi di autarchia}\\
\downarrow\\
\text{domanda di importazioni e offerta di esportazioni}\\
\downarrow\\
\text{equilibrio internazionale}\\
\downarrow\\
\text{tariffa o sussidio}\\
\downarrow\\
\text{nuovi prezzi e quantità}\\
\downarrow\\
\text{surplus e variazione del benessere}
\end{array}}
\end{equation}

Il problema economico è quindi tradotto in una successione
ordinata di operazioni che possono essere implementate in un
programma. Questa traduzione dal modello all'algoritmo costituisce
uno degli elementi centrali dell'economia computazionale.

\section{Il modello IS--LM e la fattorizzazione LU}
\label{sec:islm-lu}

Nelle sezioni precedenti abbiamo visto come un problema economico
possa essere rappresentato mediante un sistema di equazioni lineari

\begin{equation}
    Ax=b.
\end{equation}

Abbiamo inoltre introdotto l'eliminazione gaussiana, che permette
di trasformare il sistema originario in un sistema equivalente
caratterizzato da una matrice triangolare. La struttura triangolare
è particolarmente utile perché consente di determinare le incognite
in modo sequenziale.

Consideriamo ora un modello economico più articolato. Il modello
IS--LM in economia chiusa contiene più variabili endogene e ci
permette di mostrare come la soluzione di un sistema lineare possa
essere organizzata attraverso la \emph{fattorizzazione LU}.


\subsection{Il modello IS--LM}

Consideriamo il seguente modello:

\begin{align}
    Y &= C + I + \hat{G},
    \label{eq:islm-reddito}\\
    C &= \bar{C} + c(Y-\hat{T}),
    \label{eq:islm-consumo}\\
    I &= \bar{I} - bR,
    \label{eq:islm-investimento}\\
    \frac{\hat{M}}{\bar{P}} &= kY-hR.
    \label{eq:islm-moneta}
\end{align}

Le variabili endogene sono

\begin{equation}
    Y,\qquad C,\qquad I,\qquad R,
\end{equation}

dove $Y$ rappresenta il reddito, $C$ il consumo, $I$
l'investimento e $R$ il tasso di interesse.

Il modello contiene inoltre alcune variabili esogene. Seguendo la
distinzione adottata nel modello, $\bar C$, $\bar I$ e $\bar P$
sono considerate variabili esogene non controllabili, mentre

\begin{equation}
    \hat G,\qquad \hat T,\qquad \hat M
\end{equation}

sono variabili esogene controllabili dalle autorità.

Infine,

\begin{equation}
    c,\qquad b,\qquad k,\qquad h
\end{equation}

sono i parametri del modello.

Le quattro equazioni descrivono simultaneamente il funzionamento
del mercato dei beni e del mercato della moneta. Il problema
consiste nel determinare i valori di $Y$, $C$, $I$ e $R$ che
soddisfano contemporaneamente tutte le relazioni del modello.


\subsection{Dal modello economico al sistema lineare}

Per utilizzare gli strumenti numerici introdotti in precedenza,
riorganizziamo le equazioni portando le variabili endogene a
sinistra e i termini esogeni a destra.

Dall'equazione del reddito otteniamo

\begin{equation}
    Y-C-I=\hat G.
\end{equation}

Dalla funzione del consumo,

\begin{equation}
    C=\bar C+c(Y-\hat T),
\end{equation}

otteniamo

\begin{equation}
    -cY+C=\bar C-c\hat T.
\end{equation}

La funzione degli investimenti può essere scritta come

\begin{equation}
    I+bR=\bar I.
\end{equation}

Infine, l'equilibrio sul mercato monetario è

\begin{equation}
    kY-hR=\frac{\hat M}{\bar P}.
\end{equation}

Ordinando in ogni equazione le variabili secondo la sequenza
$Y,C,I,R$, il sistema diventa

\begin{align}
    Y-C-I+0R
        &=\hat G,\\
    -cY+C+0I+0R
        &=\bar C-c\hat T,\\
    0Y+0C+I+bR
        &=\bar I,\\
    kY+0C+0I-hR
        &=\frac{\hat M}{\bar P}.
\end{align}

Possiamo quindi scrivere il modello in forma matriciale:

\begin{equation}
    Ax=b,
    \label{eq:islm-matrix}
\end{equation}

dove

\begin{equation}
A=
\begin{pmatrix}
    1  & -1 & -1 & 0\\
    -c & 1  & 0  & 0\\
    0  & 0  & 1  & b\\
    k  & 0  & 0  & -h
\end{pmatrix},
\end{equation}

\begin{equation}
x=
\begin{pmatrix}
    Y\\
    C\\
    I\\
    R
\end{pmatrix},
\end{equation}

e

\begin{equation}
b=
\begin{pmatrix}
    \hat G\\
    \bar C-c\hat T\\
    \bar I\\
    \dfrac{\hat M}{\bar P}
\end{pmatrix}.
\end{equation}

Il problema economico è stato quindi ricondotto ancora una volta
alla soluzione di un sistema lineare. Rispetto al semplice modello
di domanda e offerta, il numero delle variabili è aumentato, ma la
struttura matematica del problema rimane la stessa.


\subsection{La fattorizzazione LU}

Una delle tecniche utilizzate per risolvere un sistema lineare con
matrice quadrata consiste nel fattorizzare la matrice $A$ come
prodotto di due matrici triangolari:

\begin{equation}
    \boxed{A=LU},
    \label{eq:lu-factorization}
\end{equation}

dove $L$ è una matrice triangolare inferiore
(\emph{lower triangular}) e $U$ è una matrice triangolare superiore
(\emph{upper triangular}).

La fattorizzazione LU è strettamente collegata all'idea incontrata
nell'eliminazione gaussiana. Anche in questo caso, infatti,
sfruttiamo il fatto che i sistemi caratterizzati da matrici
triangolari sono particolarmente semplici da risolvere.

La differenza è che ora la trasformazione viene rappresentata
esplicitamente attraverso la decomposizione della matrice $A$ nelle
due componenti $L$ e $U$.


\subsection{Come la fattorizzazione LU permette di risolvere il sistema}

Partiamo dal sistema

\begin{equation}
    Ax=b.
\end{equation}

Poiché $A=LU$, possiamo sostituire la fattorizzazione ottenendo

\begin{equation}
    LUx=b.
\end{equation}

Definiamo ora un vettore ausiliario

\begin{equation}
    y=Ux.
\end{equation}

Il problema può essere risolto attraverso due sistemi triangolari.

Nel primo passaggio risolviamo

\begin{equation}
    Ly=b.
    \label{eq:forward-substitution}
\end{equation}

Poiché $L$ è triangolare inferiore, il vettore $y$ può essere
determinato mediante \emph{sostituzione in avanti}.

Una volta noto $y$, risolviamo

\begin{equation}
    Ux=y.
    \label{eq:back-substitution}
\end{equation}

Poiché $U$ è triangolare superiore, possiamo determinare $x$
mediante \emph{sostituzione all'indietro}.

La procedura può quindi essere sintetizzata come

\begin{equation}
\boxed{
Ax=b
\quad\Longrightarrow\quad
LUx=b
\quad\Longrightarrow\quad
\begin{cases}
Ly=b,\\
Ux=y.
\end{cases}}
\end{equation}

Il problema originario viene così sostituito da due problemi
triangolari, che possono essere risolti in modo sequenziale.


\subsection{Esistenza della fattorizzazione LU}

La non singolarità della matrice $A$ non è, da sola, sufficiente
a garantire l'esistenza di una fattorizzazione LU nella forma
considerata.

Data una matrice

\begin{equation}
    A\in\mathbb{R}^{n\times n},
\end{equation}

la fattorizzazione LU esiste ed è unica se e solo se le
sottomatrici principali di ordine

\begin{equation}
    i=1,\ldots,n-1
\end{equation}

sono non singolari.

Le sottomatrici principali sono ottenute considerando
progressivamente le prime $i$ righe e le prime $i$ colonne della
matrice $A$.

Esistono tuttavia alcune classi di matrici per le quali l'esistenza
della fattorizzazione può essere stabilita attraverso proprietà
note. Tra queste sono particolarmente importanti le matrici a
dominanza diagonale stretta e le matrici reali simmetriche definite
positive.


\subsection{Dominanza diagonale}

Una matrice $A$ è detta \emph{a dominanza diagonale per righe} se,
per ogni riga $i$,

\begin{equation}
    |a_{ii}|
    \geq
    \sum_{j\neq i}|a_{ij}|.
    \label{eq:diagonal-dominance}
\end{equation}

In altri termini, il valore assoluto dell'elemento sulla diagonale
deve essere almeno pari alla somma dei valori assoluti degli altri
elementi appartenenti alla stessa riga.

La dominanza diagonale è detta \emph{stretta} quando

\begin{equation}
    |a_{ii}|
    >
    \sum_{j\neq i}|a_{ij}|.
    \label{eq:strict-diagonal-dominance}
\end{equation}

La stessa idea può essere applicata alle colonne, ottenendo la
definizione di dominanza diagonale per colonne.


\subsection{Matrici definite positive}

Un'altra proprietà importante è la definita positività.

Una matrice $A$ è detta \emph{definita positiva} se

\begin{equation}
    x^T A x>0
    \qquad
    \forall x\neq 0.
    \label{eq:positive-definite}
\end{equation}

L'espressione $x^TAx$ restituisce uno scalare. Infatti, se
$x\in\mathbb{R}^{n}$, allora $x^T$ ha dimensione $1\times n$,
$A$ ha dimensione $n\times n$ e $x$ ha dimensione $n\times1$.
Di conseguenza,

\begin{equation}
    x^TAx
\end{equation}

ha dimensione $1\times1$ e il suo segno può essere valutato.

Se si richiede invece

\begin{equation}
    x^TAx\geq0,
\end{equation}

la matrice viene detta \emph{semidefinita positiva}.


\subsection{Il caso particolare della fattorizzazione di Cholesky}

Quando $A$ è una matrice reale, simmetrica e definita positiva,
è possibile utilizzare una particolare fattorizzazione:

\begin{equation}
    \boxed{A=R^TR},
    \label{eq:cholesky}
\end{equation}

dove $R$ è una matrice triangolare superiore con elementi positivi
sulla diagonale.

Questa decomposizione prende il nome di
\emph{fattorizzazione di Cholesky}.

Rispetto alla fattorizzazione LU generale, in questo caso è
sufficiente determinare una sola matrice triangolare, poiché
l'altra è ottenuta mediante trasposizione. Questo riduce il numero
di operazioni necessarie.


\subsection{Come si ottengono le matrici $L$ e $U$}

Per comprendere l'idea della fattorizzazione, consideriamo per
semplicità una matrice $3\times3$.

Cerchiamo due matrici

\begin{equation}
L=
\begin{pmatrix}
    l_{11} & 0      & 0\\
    l_{21} & l_{22} & 0\\
    l_{31} & l_{32} & l_{33}
\end{pmatrix}
\end{equation}

e

\begin{equation}
U=
\begin{pmatrix}
    u_{11} & u_{12} & u_{13}\\
    0      & u_{22} & u_{23}\\
    0      & 0      & u_{33}
\end{pmatrix}
\end{equation}

tali che

\begin{equation}
    LU=A.
\end{equation}

Se lasciassimo liberi tutti gli elementi non nulli di $L$ e $U$,
avremmo più incognite che equazioni. Una possibile normalizzazione
consiste quindi nell'imporre che gli elementi della diagonale di
$L$ siano uguali a uno:

\begin{equation}
L=
\begin{pmatrix}
    1      & 0 & 0\\
    l_{21} & 1 & 0\\
    l_{31} & l_{32} & 1
\end{pmatrix}.
\end{equation}

La condizione

\begin{equation}
    LU=A
\end{equation}

genera allora un sistema di equazioni per gli elementi incogniti
di $L$ e $U$.

Ad esempio, confrontando l'elemento nella prima riga e prima
colonna otteniamo

\begin{equation}
    u_{11}=a_{11},
\end{equation}

mentre, procedendo con gli altri elementi, otteniamo
progressivamente le restanti componenti delle due matrici.

Grazie alla struttura triangolare e alla normalizzazione della
diagonale di $L$, le prime equazioni sono particolarmente semplici
da risolvere e i valori già calcolati possono essere utilizzati
nelle equazioni successive.

Nelle applicazioni non sarà tuttavia necessario costruire
manualmente $L$ e $U$: utilizzeremo le procedure di
fattorizzazione disponibili nei software di calcolo scientifico.


\subsection{Il ruolo del pivoting}

La fattorizzazione LU nella forma appena descritta può incontrare
difficoltà quando durante il procedimento gli elementi utilizzati
come \emph{pivot} non sono adatti all'eliminazione.

In questi casi è possibile modificare l'ordine delle righe della
matrice. Il \emph{pivoting} consiste nello scambiare opportunamente
le righe scegliendo elementi più adatti come pivot.

Questa estensione consente di applicare la fattorizzazione anche
in situazioni nelle quali le condizioni precedentemente illustrate
non sono direttamente soddisfatte.

Per il momento ci limitiamo a richiamare l'esistenza di questa
tecnica, rinviando a trattazioni più approfondite per la
costruzione dettagliata dell'algoritmo.


\subsection{Un ulteriore vantaggio della fattorizzazione LU}

La fattorizzazione LU può essere utile anche per calcolare il
determinante di una matrice.

Ricordiamo che

\begin{equation}
    A=LU.
\end{equation}

Poiché il determinante del prodotto di due matrici soddisfa

\begin{equation}
    \det(LU)=\det(L)\det(U),
\end{equation}

abbiamo

\begin{equation}
    \det(A)=\det(L)\det(U).
\end{equation}

Nel caso in cui la diagonale di $L$ sia costituita interamente da
elementi uguali a uno,

\begin{equation}
    \det(L)=1.
\end{equation}

Inoltre, poiché $U$ è triangolare, il suo determinante è il prodotto
degli elementi della diagonale:

\begin{equation}
    \det(U)=\prod_{i=1}^{n}u_{ii}.
\end{equation}

Di conseguenza,

\begin{equation}
    \boxed{
    \det(A)=\prod_{i=1}^{n}u_{ii}
    }.
    \label{eq:det-lu}
\end{equation}

Una volta ottenuta la fattorizzazione LU, il determinante può quindi
essere calcolato molto agevolmente, senza ricorrere allo sviluppo
di Laplace discusso in precedenza.


\subsection{Implementazione della procedura}

I principali ambienti di calcolo scientifico mettono a disposizione
funzioni specifiche per ottenere la fattorizzazione LU.

In \textsf{R}, utilizzando il pacchetto \textsf{Matrix}, possiamo
utilizzare

\begin{verbatim}
lu(A)
\end{verbatim}

In \textsf{Octave}:

\begin{verbatim}
lu(A)
\end{verbatim}

In \textsf{Python}, attraverso \textsf{SciPy}:

\begin{verbatim}
scipy.linalg.lu(A)
\end{verbatim}

mentre in \textsf{Julia} è disponibile la funzione

\begin{verbatim}
LinearAlgebra.lu(A)
\end{verbatim}

Indipendentemente dal linguaggio utilizzato, la logica
computazionale può essere sintetizzata nei seguenti passaggi:

\begin{enumerate}
    \item assegnare la matrice $A$ e il vettore $b$;
    \item ottenere la fattorizzazione di $A$ nelle matrici
          triangolari $L$ e $U$;
    \item risolvere il sistema
          \begin{equation}
              Ly=b
          \end{equation}
          per determinare $y$;
    \item risolvere il sistema
          \begin{equation}
              Ux=y
          \end{equation}
          per determinare $x$.
\end{enumerate}

Applicata al modello IS--LM, la soluzione finale

\begin{equation}
x=
\begin{pmatrix}
    Y\\
    C\\
    I\\
    R
\end{pmatrix}
\end{equation}

fornisce simultaneamente i valori di equilibrio delle quattro
variabili endogene.

L'esempio mostra quindi ancora una volta il percorso che
caratterizza l'approccio computazionale adottato in questo corso:

\begin{equation}
\boxed{
\begin{array}{c}
\text{problema economico}\\
\downarrow\\
\text{modello IS--LM}\\
\downarrow\\
\text{sistema }Ax=b\\
\downarrow\\
\text{fattorizzazione }A=LU\\
\downarrow\\
Ly=b,\quad Ux=y\\
\downarrow\\
\text{equilibrio economico}
\end{array}}
\end{equation}

\chapter{Metodi iterativi}

\section{Il modello IS--LM in economia aperta e i metodi iterativi}
\label{sec:islm-aperta-iterativi}

Nella Sezione~\ref{sec:islm-lu} abbiamo utilizzato il modello
IS--LM in economia chiusa per introdurre la fattorizzazione LU.
La fattorizzazione LU appartiene alla famiglia dei
\emph{metodi diretti}: una volta assegnati la matrice $A$ e il
vettore $b$, la soluzione del sistema

\begin{equation}
    Ax=b
\end{equation}

viene ottenuta attraverso una sequenza finita e predeterminata di
operazioni.

Consideriamo ora un'estensione del modello IS--LM nella quale due
economie sono collegate attraverso gli scambi internazionali.
L'aumento del numero delle equazioni e delle variabili ci offre
l'occasione per introdurre una diversa famiglia di tecniche
numeriche: i \emph{metodi iterativi}.

A differenza dei metodi diretti, i metodi iterativi partono da una
soluzione iniziale di prova e costruiscono una successione di
approssimazioni della soluzione. Questi metodi risultano
particolarmente interessanti quando si devono affrontare sistemi
caratterizzati da matrici grandi e sparse, cioè matrici nelle quali
molti elementi sono nulli.


\subsection{Il modello IS--LM in economia aperta}

Consideriamo due paesi. Le variabili relative al secondo paese
sono indicate con il pedice $f$.

Per il primo paese abbiamo

\begin{align}
    Y &= C+I+G+NX,\\
    C &= \bar C+c(Y-T),\\
    I &= \bar I-bR,\\
    NX &= \overline{NX}-j(Y-Y_f)
          +lE\frac{\bar P_f}{P},\\
    \frac{M}{P} &= kY-hR.
\end{align}

Per il secondo paese abbiamo invece

\begin{align}
    Y_f &= C_f+I_f+G_f+NX_f,\\
    C_f &= \bar C_f+c_f(Y_f-T_f),\\
    I_f &= \bar I_f-b_fR_f,\\
    NX_f &= -NX,\\
    \frac{M_f}{P_f} &= k_fY_f-h_fR_f.
\end{align}

Rispetto al modello in economia chiusa compare ora un nuovo
elemento: le esportazioni nette $NX$ collegano le due economie.

La relazione

\begin{equation}
    NX_f=-NX
\end{equation}

esprime il fatto che le esportazioni nette di un paese
corrispondono alle importazioni nette dell'altro.

Date le variabili esogene e i parametri, il problema consiste
ancora una volta nel determinare simultaneamente le variabili
endogene del modello. Una volta ottenuto l'equilibrio, possiamo
inoltre modificare gli strumenti di politica fiscale o monetaria
e studiarne gli effetti sulle due economie.

Anche questo modello può essere ricondotto alla forma

\begin{equation}
    Ax=b.
\end{equation}

Il problema matematico rimane quindi quello già incontrato nelle
sezioni precedenti. Cambia, tuttavia, la tecnica numerica che
vogliamo utilizzare per determinarne la soluzione.


\subsection{L'idea dei metodi iterativi}

Un metodo iterativo non cerca necessariamente di ottenere la
soluzione attraverso una sequenza prefissata di operazioni.
Partendo da un vettore iniziale $x_0$, genera invece una successione

\begin{equation}
    x_0,\;x_1,\;x_2,\;\ldots
\end{equation}

che, se il metodo converge, si avvicina progressivamente alla
soluzione del sistema.

Per costruire un meccanismo iterativo partiamo ancora da

\begin{equation}
    Ax=b.
    \label{eq:iter-system}
\end{equation}

Introduciamo una matrice $M$ e scriviamo

\begin{equation}
    A=M-(M-A).
\end{equation}

Sostituendo questa espressione nel sistema otteniamo

\begin{equation}
    \left[M-(M-A)\right]x=b,
\end{equation}

da cui

\begin{equation}
    Mx=(M-A)x+b.
\end{equation}

Introducendo gli indici di iterazione, possiamo utilizzare i valori
calcolati all'iterazione $k$ per costruire una nuova approssimazione:

\begin{equation}
    Mx_{k+1}=(M-A)x_k+b.
\end{equation}

Sottraendo $Mx_k$ da entrambi i membri otteniamo

\begin{equation}
    M(x_{k+1}-x_k)=b-Ax_k.
    \label{eq:iter-general}
\end{equation}

Questa relazione costituisce la base del meccanismo iterativo.


\subsection{Correzione e residuo}

Per interpretare l'Equazione~\eqref{eq:iter-general}, definiamo

\begin{equation}
    z_k=x_{k+1}-x_k
    \label{eq:correction}
\end{equation}

e

\begin{equation}
    r_k=b-Ax_k.
    \label{eq:residual}
\end{equation}

Il vettore $z_k$ rappresenta la \emph{correzione} che deve essere
aggiunta all'approssimazione corrente per ottenere quella
successiva.

Il vettore $r_k$ prende invece il nome di \emph{residuo}. Esso
misura quanto l'approssimazione corrente $x_k$ non soddisfa ancora
il sistema originale.

Se $x_k$ fosse già la soluzione esatta, avremmo infatti

\begin{equation}
    Ax_k=b
\end{equation}

e quindi

\begin{equation}
    r_k=0.
\end{equation}

Utilizzando queste definizioni, l'Equazione~\eqref{eq:iter-general}
diventa

\begin{equation}
    Mz_k=r_k.
    \label{eq:Mzr}
\end{equation}

Il problema iterativo può quindi essere interpretato in modo molto
semplice: calcoliamo l'errore commesso dall'approssimazione corrente,
espresso dal residuo, e utilizziamo questa informazione per
determinare una correzione $z_k$.


\subsection{Il meccanismo iterativo generale}

Possiamo ora descrivere la procedura completa.

Si parte scegliendo arbitrariamente un vettore iniziale

\begin{equation}
    x_0.
\end{equation}

Si calcola il residuo

\begin{equation}
    r_0=b-Ax_0.
\end{equation}

Successivamente si risolve il sistema

\begin{equation}
    Mz_0=r_0
\end{equation}

per determinare la correzione $z_0$.

La nuova approssimazione è quindi

\begin{equation}
    x_1=x_0+z_0.
\end{equation}

La stessa procedura viene ripetuta:

\begin{equation}
    r_1=b-Ax_1,
\end{equation}

\begin{equation}
    Mz_1=r_1,
\end{equation}

\begin{equation}
    x_2=x_1+z_1,
\end{equation}

e così via.

In termini generali, a ogni iterazione $k$ eseguiamo quindi i tre
passaggi

\begin{align}
    r_k &= b-Ax_k,\\
    Mz_k &= r_k,\\
    x_{k+1} &= x_k+z_k.
\end{align}

Questo schema costituisce il meccanismo iterativo generale
presentato nelle slide come \emph{metodo di Richardson}.


\subsection{La scelta della matrice $M$}

A ogni iterazione dobbiamo risolvere il sistema

\begin{equation}
    Mz_k=r_k.
\end{equation}

Questa operazione rappresenta la parte computazionalmente più
impegnativa del procedimento.

Se la soluzione di questo sistema fosse difficile quanto quella
del sistema originale $Ax=b$, il metodo iterativo non offrirebbe
un particolare vantaggio. La matrice $M$ deve quindi essere scelta
in modo tale che il sistema $Mz_k=r_k$ sia molto semplice da
risolvere.

Due scelte particolarmente convenienti sono una matrice diagonale
oppure una matrice triangolare.

Scelte differenti di $M$ danno origine a differenti metodi
iterativi.

Nel \emph{metodo di Jacobi} si utilizza

\begin{equation}
    M=\operatorname{diag}(A),
\end{equation}

cioè la matrice contenente gli elementi della diagonale di $A$ e
zeri in tutte le altre posizioni.

Nel \emph{metodo di Gauss--Seidel}, invece, $M$ è costituita dalla
parte triangolare inferiore di $A$.

Poiché una matrice diagonale o triangolare rende particolarmente
semplice la soluzione di $Mz_k=r_k$, ciascuna iterazione può essere
eseguita con un costo relativamente contenuto.

Le slide introducono inoltre il \emph{metodo del rilassamento},
nel quale la diagonale di $M$ viene modificata attraverso un
parametro $\omega\neq0$.


\subsection{Convergenza dei metodi stazionari}

La semplicità di una singola iterazione non garantisce che la
successione

\begin{equation}
    x_0,x_1,x_2,\ldots
\end{equation}

converga effettivamente alla soluzione.

Le condizioni di convergenza dipendono dal metodo e dalle
caratteristiche della matrice $A$.

Per il metodo di Jacobi, una condizione indicata nelle slide è
che $A$ sia a dominanza diagonale stretta per righe.

Per Gauss--Seidel, le condizioni considerate sono:

\begin{itemize}
    \item $A$ a dominanza diagonale stretta per righe;
    \item $A$ simmetrica e definita positiva.
\end{itemize}

Queste restrizioni rappresentano un limite importante dei metodi
iterativi considerati fino a questo punto.


\subsection{Introduzione di un parametro di passo}

Il meccanismo precedente aggiorna la soluzione secondo la regola

\begin{equation}
    x_{k+1}=x_k+z_k.
\end{equation}

Una possibile estensione consiste nell'introdurre un parametro
$\alpha_k$:

\begin{equation}
    \boxed{
    x_{k+1}=x_k+\alpha_k z_k
    }.
    \label{eq:iter-alpha}
\end{equation}

Il termine $\alpha_kz_k$ determina quindi il movimento effettuato
a ogni iterazione.

Possiamo interpretare separatamente le due componenti:

\begin{itemize}
    \item $z_k$ determina la \emph{direzione} del movimento;
    \item $\alpha_k$ determina l'\emph{intensità} del movimento.
\end{itemize}

Per ottenere una convergenza rapida è quindi importante scegliere
in modo appropriato sia la direzione sia l'ampiezza del passo.


\subsection{Metodi stazionari e non stazionari}

La presenza del parametro $\alpha_k$ permette di introdurre una
distinzione tra due famiglie di metodi iterativi.

Nei \emph{metodi stazionari}, il parametro non cambia nel corso
delle iterazioni:

\begin{equation}
    \alpha_k=\alpha.
\end{equation}

I metodi di Jacobi e Gauss--Seidel appartengono alla famiglia dei
metodi stazionari.

Nei \emph{metodi non stazionari}, invece, il parametro può cambiare
a ogni iterazione:

\begin{equation}
    \alpha_k
\end{equation}

viene quindi scelto in funzione dell'informazione disponibile al
passo $k$.

I metodi stazionari hanno il vantaggio di essere semplici, ma
convergono soltanto sotto determinate condizioni. Questo motiva
l'introduzione di metodi nei quali direzione e intensità del
movimento possono essere adattate nel corso delle iterazioni.


\subsection{Il caso di una matrice simmetrica e definita positiva}

Consideriamo ora il caso in cui $A$ sia simmetrica e definita
positiva.

Associamo al sistema $Ax=b$ la funzione

\begin{equation}
    \Phi(x)
    =
    \frac{1}{2}x^TAx-x^Tb.
    \label{eq:phi-gradient}
\end{equation}

Il gradiente di questa funzione è

\begin{equation}
    \nabla\Phi(x)=Ax-b.
    \label{eq:gradient-phi}
\end{equation}

Se $A$ è simmetrica e definita positiva, $\Phi(x)$ è strettamente
convessa e possiede quindi un unico punto di minimo.

Nel punto di minimo il gradiente è nullo:

\begin{equation}
    \nabla\Phi(x)=0.
\end{equation}

Utilizzando l'Equazione~\eqref{eq:gradient-phi}, questa condizione
diventa

\begin{equation}
    Ax-b=0,
\end{equation}

ovvero

\begin{equation}
    Ax=b.
\end{equation}

La soluzione di un sistema lineare con matrice simmetrica e
definita positiva può quindi essere interpretata anche come il
punto che minimizza la funzione quadratica $\Phi(x)$.

Abbiamo così trasformato il problema di soluzione del sistema in
un problema di minimizzazione.


\subsection{Il gradiente come direzione}

Per costruire un metodo iterativo dobbiamo ora scegliere una
direzione $z_k$ che ci permetta di avvicinarci al minimo della
funzione $\Phi$.

Il gradiente

\begin{equation}
    \nabla\Phi(x)
\end{equation}

indica la direzione di massima crescita della funzione. Per
comprendere questa proprietà consideriamo la derivata direzionale
di una generica funzione $f$ nel punto $x_0$, lungo una direzione
rappresentata dal versore $u$:

\begin{equation}
    D_uf(x_0)=\nabla f(x_0)\cdot u,
    \qquad
    \|u\|=1.
\end{equation}

Utilizzando la definizione di prodotto scalare possiamo scrivere

\begin{equation}
    \nabla f(x_0)\cdot u
    =
    \|\nabla f(x_0)\|\,\|u\|\cos\theta,
\end{equation}

dove $\theta$ è l'angolo tra $\nabla f(x_0)$ e $u$.

Poiché $u$ è un versore,

\begin{equation}
    \|u\|=1,
\end{equation}

e quindi

\begin{equation}
    D_uf(x_0)
    =
    \|\nabla f(x_0)\|\cos\theta.
\end{equation}

Il valore massimo si ottiene quando

\begin{equation}
    \cos\theta=1,
\end{equation}

cioè quando $\theta=0$. In questo caso $u$ ha la stessa direzione
del gradiente:

\begin{equation}
    u=
    \frac{\nabla f(x_0)}
         {\|\nabla f(x_0)\|}.
\end{equation}

La massima derivata direzionale è pertanto

\begin{equation}
    \max_{\|u\|=1}D_uf(x_0)
    =
    \|\nabla f(x_0)\|.
\end{equation}

Il gradiente indica dunque la direzione di massima crescita.
Se vogliamo invece ridurre $\Phi$, dobbiamo muoverci nel verso
opposto.

Una possibile scelta della direzione è quindi

\begin{equation}
    z_k=-\nabla\Phi(x_k).
    \label{eq:steepest-direction}
\end{equation}


\subsection{La scelta dell'intensità del movimento}

Una volta individuata la direzione $z_k$, rimane da stabilire
quanto muoversi lungo quella direzione.

Partendo dalla regola

\begin{equation}
    x_{k+1}=x_k+\alpha_kz_k,
\end{equation}

possiamo scegliere $\alpha_k$ in modo da minimizzare la funzione
lungo la direzione $z_k$:

\begin{equation}
    \alpha_k
    =
    \arg\min_{\alpha}
    \Phi(x_k+\alpha z_k).
    \label{eq:line-search}
\end{equation}

La combinazione

\begin{equation}
    z_k=-\nabla\Phi(x_k)
\end{equation}

con la scelta dell'ampiezza del passo conduce al
\emph{metodo del gradiente}.

Quando $A$ è simmetrica e definita positiva, il metodo converge
alla soluzione a partire da qualsiasi punto iniziale $x_0$.

Un'evoluzione di questa idea è il \emph{metodo del gradiente
coniugato}, nel quale vengono costruite direzioni che permettono
di accelerare la convergenza.


\subsection{Il caso di una matrice generica}

La situazione è più complessa quando $A$ non è simmetrica e
definita positiva.

La soluzione iterativa di sistemi con matrici generiche costituisce
un argomento molto ampio. In questo corso non ne svilupperemo in
dettaglio la teoria, ma utilizzeremo uno dei metodi maggiormente
diffusi: il \emph{Generalized Minimal Residual Method}, indicato
con l'acronimo GMRES.

Il metodo GMRES appartiene alla famiglia dei metodi che costruiscono
le nuove approssimazioni cercando di ridurre il residuo.

Nelle applicazioni utilizzeremo direttamente le implementazioni
disponibili nei principali software scientifici. Ad esempio:

\begin{verbatim}
# R
pracma::gmres(A, b, ...)

# Octave
gmres(A, b, ...)

# Python
scipy.sparse.linalg.gmres(A, b, ...)

# Julia
IterativeSolvers.gmres(A, b, ...)
\end{verbatim}

L'obiettivo, in questa fase, non è ricostruire l'algoritmo GMRES,
ma comprendere perché possa essere necessario ricorrere a un metodo
iterativo più generale quando le proprietà della matrice non
permettono di utilizzare convenientemente i metodi precedenti.


\subsection{Una classificazione dei metodi per sistemi lineari}

Possiamo a questo punto riordinare le tecniche introdotte finora
in due grandi famiglie.

I \emph{metodi diretti} determinano la soluzione attraverso un
numero predefinito di operazioni. In questa categoria abbiamo
incontrato:

\begin{itemize}
    \item l'eliminazione di Gauss;
    \item la fattorizzazione LU.
\end{itemize}

I \emph{metodi iterativi}, invece, costruiscono una successione
di approssimazioni della soluzione. Possiamo distinguerli
ulteriormente in metodi stazionari e non stazionari.

Tra i metodi stazionari abbiamo considerato:

\begin{itemize}
    \item lo schema generale di Richardson;
    \item il metodo di Jacobi;
    \item il metodo di Gauss--Seidel;
    \item il metodo del rilassamento.
\end{itemize}

Tra i metodi non stazionari, la scelta dipende anche dalle
caratteristiche della matrice. Per matrici simmetriche e definite
positive abbiamo introdotto:

\begin{itemize}
    \item il metodo del gradiente;
    \item il metodo del gradiente coniugato.
\end{itemize}

Per matrici generiche abbiamo invece introdotto il metodo

\begin{equation}
    \text{GMRES}.
\end{equation}

La scelta del metodo numerico non è quindi indipendente dalla
struttura del problema. Le dimensioni della matrice, la presenza
di molti elementi nulli e proprietà quali simmetria, definita
positività e dominanza diagonale possono rendere un metodo più
appropriato di un altro.


\subsection{Il ruolo del residuo nei metodi iterativi}

Un elemento comune ai metodi iterativi è il residuo

\begin{equation}
    r_k=b-Ax_k.
\end{equation}

Esso fornisce informazioni sulla qualità dell'approssimazione
corrente e può essere utilizzato sia per determinare come
procedere verso una nuova approssimazione, sia per stabilire
quando interrompere il procedimento.

I metodi iterativi più avanzati costruiscono le nuove
approssimazioni all'interno di particolari sottospazi, detti
\emph{sottospazi di Krylov}, che si ampliano progressivamente
con il procedere delle iterazioni.

All'interno di questa famiglia possono essere distinti due
approcci richiamati nelle slide.

Nel primo, detto \emph{Ritz--Galerkin} o del residuo ortogonale,
la nuova approssimazione viene scelta imponendo una condizione
di ortogonalità sul residuo. A questa famiglia appartengono il
gradiente coniugato, utilizzato per matrici simmetriche e definite
positive, e il bi-CG per matrici generiche.

Nel secondo, detto \emph{approccio del minimo residuo}, si sceglie
invece l'approssimazione che rende minima l'ampiezza del residuo
all'interno del sottospazio considerato. Il metodo GMRES appartiene
a questa seconda famiglia.

Questi concetti forniscono un'indicazione della logica sulla quale
si basano i metodi iterativi moderni. In questo corso non
svilupperemo ulteriormente la teoria dei sottospazi di Krylov, ma
utilizzeremo gli algoritmi disponibili nei software di calcolo
scientifico quando le applicazioni lo richiederanno.


\subsection{Dal problema economico alla soluzione iterativa}

Possiamo infine sintetizzare il percorso seguito in questa sezione:

\begin{equation}
\boxed{
\begin{array}{c}
\text{modello IS--LM in economia aperta}\\
\downarrow\\
\text{sistema lineare }Ax=b\\
\downarrow\\
\text{scelta di un'approssimazione iniziale }x_0\\
\downarrow\\
r_k=b-Ax_k\\
\downarrow\\
\text{determinazione della correzione}\\
\downarrow\\
x_{k+1}=x_k+\alpha_kz_k\\
\downarrow\\
\text{nuova approssimazione della soluzione}
\end{array}}
\end{equation}

Il passaggio dai metodi diretti ai metodi iterativi introduce
quindi un modo diverso di affrontare il problema $Ax=b$.
Invece di costruire direttamente la soluzione, si genera una
sequenza di approssimazioni che viene progressivamente corretta
utilizzando l'informazione contenuta nel residuo.

\section{Implementazione dei metodi iterativi}
\label{sec:implementazione-metodi-iterativi}

Nella sezione precedente abbiamo introdotto i metodi iterativi per
la soluzione di un sistema lineare

\begin{equation}
    Ax=b.
\end{equation}

L'idea fondamentale consiste nel partire da un'approssimazione
iniziale $x_0$ e costruire una successione

\begin{equation}
    x_0,\;x_1,\;x_2,\;\ldots
\end{equation}

che, sotto opportune condizioni, converge alla soluzione del sistema.

In termini generali, a ogni iterazione viene calcolato il residuo

\begin{equation}
    r_k=b-Ax_k,
\end{equation}

dal quale si ricava una correzione della soluzione corrente.
Il procedimento viene quindi ripetuto fino a quando
l'approssimazione ottenuta può essere considerata sufficientemente
accurata.

Dal punto di vista teorico, il meccanismo è relativamente semplice.
La sua implementazione pone però due nuovi problemi computazionali:

\begin{enumerate}
    \item come possiamo chiedere al computer di ripetere più volte
          la stessa sequenza di operazioni?
    \item come possiamo stabilire quando il procedimento iterativo
          deve essere interrotto?
\end{enumerate}

Per rispondere alla prima domanda introduciamo i cicli
\texttt{for} e \texttt{while}. Per rispondere alla seconda
introduciamo il concetto di norma e lo utilizziamo per misurare
la dimensione del residuo.


\subsection{Ripetere un'operazione: il ciclo \texttt{for}}

Un algoritmo iterativo richiede di eseguire ripetutamente la stessa
sequenza di operazioni.

Supponiamo, ad esempio, di voler eseguire un'operazione per un numero
prefissato di volte. In questo caso possiamo utilizzare un ciclo
\texttt{for}.

La struttura generale in \textsf{R} è

\begin{verbatim}
for (i in sequenza) {

    operazioni

}
\end{verbatim}

La variabile $i$ prende, uno dopo l'altro, i valori contenuti nella
sequenza. Per ciascun valore vengono eseguite le istruzioni
contenute nel ciclo.

Ad esempio,

\begin{verbatim}
for (i in 1:10) {

    print(i)

}
\end{verbatim}

esegue il comando \texttt{print(i)} dieci volte, assegnando
progressivamente a \texttt{i} i valori da 1 a 10.

Nel contesto di un metodo iterativo, l'indice $i$ può essere
interpretato come l'indice dell'iterazione:

\begin{equation}
    i=1,2,\ldots,K,
\end{equation}

dove $K$ rappresenta il numero massimo di iterazioni che vogliamo
eseguire.

Una struttura di questo tipo permette quindi di tradurre in codice
una successione del tipo

\begin{equation}
    x_0
    \longrightarrow
    x_1
    \longrightarrow
    x_2
    \longrightarrow
    \cdots
    \longrightarrow
    x_K.
\end{equation}


\subsection{Ripetere un'operazione finché una condizione è soddisfatta}

Nei metodi iterativi, tuttavia, non sempre sappiamo in anticipo
quante iterazioni saranno necessarie.

Potremmo voler continuare il procedimento fino a quando la soluzione
raggiunge un determinato livello di accuratezza.

In questo caso è più naturale utilizzare un ciclo \texttt{while}.

La struttura generale è

\begin{verbatim}
while (condizione) {

    operazioni

}
\end{verbatim}

Le istruzioni contenute nel ciclo vengono eseguite finché la
condizione specificata rimane vera.

Consideriamo, ad esempio,

\begin{verbatim}
i <- 1

while (i < 10) {

    print(i)
    i <- i + 1

}
\end{verbatim}

All'inizio abbiamo $i=1$. Dopo ogni iterazione il valore di $i$
viene aumentato di una unità. Il ciclo continua finché

\begin{equation}
    i<10.
\end{equation}

Quando questa condizione non è più verificata, il ciclo termina.

Il ciclo \texttt{while} è particolarmente adatto ai metodi
iterativi perché permette di formulare una regola del tipo:

\begin{quote}
continua a costruire nuove approssimazioni finché l'errore è
troppo grande.
\end{quote}

Per rendere operativa questa idea dobbiamo però definire come
misurare l'errore.


\subsection{Distanza e norma}

Per misurare quanto due numeri reali $a$ e $b$ siano diversi
possiamo utilizzare la distanza

\begin{equation}
    d(a,b)=|a-b|.
\end{equation}

Quando lavoriamo con vettori, la stessa idea richiede una
generalizzazione. Questa generalizzazione è fornita dal concetto
di \emph{norma}.

Una norma associa a un vettore un numero reale non negativo che
può essere interpretato come una misura della sua grandezza.

Per un vettore

\begin{equation}
    x=
    \begin{pmatrix}
        x_1\\
        x_2\\
        \vdots\\
        x_n
    \end{pmatrix},
\end{equation}

possiamo definire diverse norme.


\subsubsection{Norma 1}

La norma 1 è definita come

\begin{equation}
    \|x\|_1
    =
    \sum_{i=1}^{n}|x_i|.
    \label{eq:norma1}
\end{equation}

Essa corrisponde quindi alla somma dei valori assoluti degli
elementi del vettore.


\subsubsection{Norma 2}

La norma 2, detta anche norma euclidea, è

\begin{equation}
    \|x\|_2
    =
    \sqrt{
        \sum_{i=1}^{n}x_i^2
    }.
    \label{eq:norma2}
\end{equation}

Nel caso di un vettore con due componenti,

\begin{equation}
    x=
    \begin{pmatrix}
        x_1\\
        x_2
    \end{pmatrix},
\end{equation}

otteniamo

\begin{equation}
    \|x\|_2
    =
    \sqrt{x_1^2+x_2^2}.
\end{equation}

Questa espressione coincide con la distanza euclidea del punto
$(x_1,x_2)$ dall'origine.


\subsubsection{Norma infinito}

La norma infinito è invece definita come

\begin{equation}
    \|x\|_\infty
    =
    \max_i |x_i|.
    \label{eq:normainf}
\end{equation}

Essa misura quindi la componente del vettore che presenta il
valore assoluto maggiore.

Le tre norme forniscono misure differenti della grandezza di un
vettore. Nel seguito utilizzeremo principalmente la norma 2.


\subsection{La norma in \textsf{R}}

In \textsf{R}, la funzione \texttt{norm()} opera su oggetti di
tipo matrice. Se $x$ è inizialmente un vettore, possiamo quindi
trasformarlo in una matrice mediante

\begin{verbatim}
x <- as.matrix(x)
\end{verbatim}

e successivamente calcolarne la norma.

La norma 1 può essere ottenuta con

\begin{verbatim}
norm(x, type = "1")
\end{verbatim}

la norma 2 con

\begin{verbatim}
norm(x, type = "2")
\end{verbatim}

e la norma infinito con

\begin{verbatim}
norm(x, type = "I")
\end{verbatim}

La possibilità di calcolare numericamente la norma di un vettore
ci permette ora di costruire un criterio per decidere quando
interrompere un algoritmo iterativo.


\subsection{Il residuo come misura della qualità dell'approssimazione}

Riprendiamo il sistema

\begin{equation}
    Ax=b.
\end{equation}

Se $x_k$ è l'approssimazione disponibile all'iterazione $k$,
definiamo il residuo come

\begin{equation}
    r_k=b-Ax_k.
    \label{eq:residuo-implementazione}
\end{equation}

Se $x_k$ fosse la soluzione esatta, avremmo

\begin{equation}
    Ax_k=b
\end{equation}

e pertanto

\begin{equation}
    r_k=0.
\end{equation}

Se invece $x_k$ è soltanto un'approssimazione, il residuo sarà in
generale diverso dal vettore nullo.

Possiamo quindi utilizzare la norma del residuo,

\begin{equation}
    \|r_k\|,
\end{equation}

come misura di quanto l'approssimazione corrente non soddisfa
ancora il sistema.

Più precisamente, possiamo considerare

\begin{equation}
    \|r_k\|_2
    =
    \|b-Ax_k\|_2.
    \label{eq:norma-residuo}
\end{equation}

Se il metodo converge, ci aspettiamo che questa quantità diventi
progressivamente più piccola.


\subsection{Tolleranza e criterio di arresto}

In un calcolo numerico non è generalmente necessario richiedere
che il residuo sia esattamente uguale a zero.

Definiamo invece una quantità positiva molto piccola,
che indichiamo con $\varepsilon$ e chiamiamo \emph{tolleranza}.

Possiamo considerare l'approssimazione sufficientemente accurata
quando

\begin{equation}
    \|r_k\|_2
    \leq
    \varepsilon.
    \label{eq:stopping-rule}
\end{equation}

Il procedimento deve quindi continuare mentre

\begin{equation}
    \|r_k\|_2
    >
    \varepsilon.
\end{equation}

Questa condizione si traduce naturalmente in un ciclo
\texttt{while}:

\begin{verbatim}
while (norm(r, type = "2") > tolleranza) {

    # nuova iterazione

}
\end{verbatim}

La tolleranza stabilisce quindi il livello di precisione richiesto.
Una tolleranza più piccola impone un criterio di arresto più
stringente e può richiedere un numero maggiore di iterazioni.


\subsection{Costruzione dell'algoritmo iterativo}

Possiamo ora combinare gli strumenti introdotti in questa sezione
con il meccanismo iterativo studiato nella sezione precedente.

Partiamo dal sistema

\begin{equation}
    Ax=b.
\end{equation}

Assegniamo un'approssimazione iniziale

\begin{equation}
    x_0
\end{equation}

e calcoliamo il residuo iniziale

\begin{equation}
    r_0=b-Ax_0.
\end{equation}

Finché la norma del residuo è superiore alla tolleranza,
continuiamo a calcolare nuove approssimazioni.

Lo schema generale dell'algoritmo può essere rappresentato come

\begin{verbatim}
scegli x0
calcola r = b - A x0

while (norma(r) > tolleranza) {

    calcola la correzione z
    aggiorna x
    calcola il nuovo residuo r

}
\end{verbatim}

In termini matematici, utilizzando il meccanismo introdotto
precedentemente, a ogni iterazione calcoliamo

\begin{equation}
    r_k=b-Ax_k,
\end{equation}

risolviamo

\begin{equation}
    Mz_k=r_k
\end{equation}

e aggiorniamo la soluzione secondo

\begin{equation}
    x_{k+1}=x_k+z_k.
\end{equation}

Il procedimento continua fino a quando

\begin{equation}
    \|b-Ax_k\|_2
    \leq
    \varepsilon.
\end{equation}


\subsection{Il rischio di un ciclo infinito}

L'utilizzo di un ciclo \texttt{while} richiede una particolare
attenzione.

Se il metodo iterativo non converge, oppure se la condizione di
arresto non può essere soddisfatta, il ciclo potrebbe continuare
indefinitamente.

Per evitare questo problema è opportuno introdurre anche un numero
massimo di iterazioni, che indichiamo con $K_{\max}$.

Il procedimento viene quindi interrotto quando si verifica almeno
una delle seguenti condizioni:

\begin{equation}
    \|r_k\|_2\leq\varepsilon
\end{equation}

oppure

\begin{equation}
    k\geq K_{\max}.
\end{equation}

La struttura generale diventa quindi

\begin{verbatim}
while (norm(r, type = "2") > tolleranza &&
       iterazione < max_iterazioni) {

    # calcolo della nuova approssimazione

}
\end{verbatim}

Il numero massimo di iterazioni costituisce una misura di
sicurezza: impedisce al programma di continuare indefinitamente
nel caso in cui il metodo non converga.


\subsection{Interpretazione dell'algoritmo}

Possiamo infine rappresentare l'intera procedura come

\begin{equation}
\boxed{
\begin{array}{c}
\text{scegliere }x_0\\
\downarrow\\
r_0=b-Ax_0\\
\downarrow\\
\text{calcolare }\|r_0\|\\
\downarrow\\
\|r_k\|\leq\varepsilon\;?
\end{array}}
\end{equation}

Se la risposta è positiva, il procedimento termina e $x_k$ viene
accettato come soluzione numerica.

Se la risposta è negativa, viene calcolata una nuova
approssimazione:

\begin{equation}
\boxed{
x_k
\longrightarrow
r_k
\longrightarrow
z_k
\longrightarrow
x_{k+1}
\longrightarrow
r_{k+1}
}
\end{equation}

e il criterio di arresto viene verificato nuovamente.

Questa struttura evidenzia il ruolo dei diversi strumenti
computazionali introdotti nella sezione:

\begin{itemize}
    \item il ciclo \texttt{for} permette di ripetere una procedura
          per un numero prefissato di volte;
    \item il ciclo \texttt{while} permette di ripeterla finché
          una determinata condizione rimane verificata;
    \item la norma permette di associare a un vettore una misura
          scalare della sua grandezza;
    \item il residuo misura quanto l'approssimazione corrente non
          soddisfa ancora il sistema;
    \item la tolleranza stabilisce quando l'approssimazione può
          essere considerata sufficientemente accurata.
\end{itemize}

Gli strumenti introdotti in questa sezione permettono quindi di
passare dalla formulazione matematica di un metodo iterativo alla
sua effettiva implementazione in un linguaggio di programmazione.

\chapter{applicazione alle valutazioni di impatto ambientale}

\section{Valutazione di impatto ambientale: Life Cycle Assessment}
\label{sec:lca}

Gli strumenti dell'algebra lineare e del calcolo numerico introdotti
nelle sezioni precedenti trovano applicazione anche in ambiti che,
a prima vista, possono apparire molto distanti dai modelli economici
considerati finora.

Un esempio particolarmente importante è rappresentato dalla
\emph{Life Cycle Assessment} (LCA), una metodologia utilizzata per
valutare gli impatti ambientali associati a un prodotto, a un servizio
o, più in generale, a un sistema produttivo.

L'idea fondamentale consiste nel considerare non soltanto ciò che
avviene nell'ultima fase del processo produttivo, ma l'insieme dei
processi necessari per ottenere il prodotto desiderato. La produzione
di un bene può infatti richiedere altri beni intermedi, la cui
produzione richiede a sua volta ulteriori input. Ciascuno di questi
processi può inoltre utilizzare risorse naturali e generare emissioni
nell'ambiente.

La LCA cerca quindi di ricostruire queste relazioni e di rispondere,
in forma semplificata, a due domande:

\begin{enumerate}
    \item quali flussi ambientali sono necessari o generati per
          ottenere una determinata quantità di prodotto?
    \item quali impatti ambientali possono essere associati a tali
          flussi?
\end{enumerate}

Queste due domande corrispondono alle due fasi della LCA sulle quali
ci concentreremo: la \emph{Life Cycle Inventory} (LCI) e la
\emph{Life Cycle Impact Assessment} (LCIA).


\subsection{Le fasi della Life Cycle Assessment}

Secondo lo standard ISO 14040, una Life Cycle Assessment comprende
quattro fasi:

\begin{enumerate}
    \item definizione dell'obiettivo e del campo di applicazione
          (\emph{goal and scope definition});
    \item analisi dell'inventario
          (\emph{Life Cycle Inventory}, LCI);
    \item valutazione degli impatti
          (\emph{Life Cycle Impact Assessment}, LCIA);
    \item interpretazione dei risultati.
\end{enumerate}

La prima fase stabilisce, tra gli altri aspetti, quale sistema viene
studiato e quale risultato deve essere ottenuto dall'analisi.
L'ultima fase riguarda invece l'interpretazione dei risultati
ottenuti.

In questa sezione ci concentreremo sulle due fasi centrali, LCI e
LCIA, perché permettono di mettere particolarmente bene in evidenza
la struttura matematica e computazionale della LCA.

Il percorso che seguiremo può essere anticipato schematicamente come

\begin{equation}
\boxed{
\text{domanda finale}
\longrightarrow
\text{flussi ambientali}
\longrightarrow
\text{impatti ambientali}
}.
\end{equation}


\subsection{La fase di inventario: Life Cycle Inventory}

La fase LCI descrive il sistema produttivo attraverso i suoi input
e output.

Un \emph{sistema di prodotto} è composto da diversi
\emph{processi unitari}. Ogni processo può:

\begin{itemize}
    \item utilizzare beni prodotti da altri processi;
    \item produrre beni o servizi;
    \item utilizzare risorse prelevate dall'ambiente;
    \item generare emissioni nell'ambiente.
\end{itemize}

Per comprendere la struttura computazionale della LCI utilizziamo
un esempio estremamente semplice, costituito inizialmente da un
solo processo: la produzione di elettricità.


\subsection{Rappresentazione di un processo produttivo}

Supponiamo che un processo produca $10$ kWh di elettricità
utilizzando combustibile e generando emissioni di anidride carbonica
e biossido di zolfo.

Per rappresentare il processo dobbiamo innanzitutto identificare
tutti gli elementi, o \emph{item}, coinvolti.

Introduciamo quindi il vettore

\begin{equation}
p_i=
\begin{pmatrix}
\text{litri di combustibile}\\
\text{kWh di elettricità}\\
\text{kg di anidride carbonica}\\
\text{kg di biossido di zolfo}
\end{pmatrix}.
\end{equation}

Il vettore $p_i$ non contiene quantità, ma identifica gli elementi
del sistema e le rispettive unità di misura. Esso costituisce quindi
una struttura di riferimento che permette di attribuire un
significato economico o ambientale alle diverse posizioni dei
vettori numerici che utilizzeremo.

Il processo di produzione dell'elettricità può essere rappresentato
dal vettore

\begin{equation}
p_1=
\begin{pmatrix}
-2\\
10\\
1\\
0.1
\end{pmatrix}.
\end{equation}

La posizione degli elementi di $p_1$ corrisponde esattamente alla
posizione degli elementi del vettore $p_i$. Possiamo quindi leggere
il vettore nel modo seguente:

\begin{itemize}
    \item $-2$: il processo utilizza $2$ litri di combustibile;
    \item $10$: il processo produce $10$ kWh di elettricità;
    \item $1$: il processo emette $1$ kg di $\mathrm{CO_2}$;
    \item $0.1$: il processo emette $0.1$ kg di $\mathrm{SO_2}$.
\end{itemize}

Adottiamo quindi una convenzione sui segni: un valore positivo
indica un output del processo, mentre un valore negativo indica un
input utilizzato dal processo.


\subsection{Aggiunta di un secondo processo}

Il processo di produzione dell'elettricità utilizza combustibile.
Se vogliamo considerare anche gli effetti associati alla produzione
di questo combustibile, dobbiamo aggiungere un secondo processo al
sistema.

Supponiamo che il processo di produzione del combustibile utilizzi
petrolio greggio. Compare quindi un nuovo elemento e il vettore
degli item deve essere aggiornato:

\begin{equation}
p_i=
\begin{pmatrix}
\text{litri di combustibile}\\
\text{kWh di elettricità}\\
\text{kg di anidride carbonica}\\
\text{kg di biossido di zolfo}\\
\text{litri di petrolio greggio}
\end{pmatrix}.
\end{equation}

Anche il vettore del primo processo deve essere aggiornato
aggiungendo uno zero nella posizione corrispondente al petrolio
greggio:

\begin{equation}
p_1=
\begin{pmatrix}
-2\\
10\\
1\\
0.1\\
0
\end{pmatrix}.
\end{equation}

Supponiamo che il secondo processo sia rappresentato da

\begin{equation}
p_2=
\begin{pmatrix}
100\\
0\\
10\\
2\\
-50
\end{pmatrix}.
\end{equation}

Ciò significa che una attivazione del processo:

\begin{itemize}
    \item produce $100$ litri di combustibile;
    \item non produce elettricità;
    \item emette $10$ kg di $\mathrm{CO_2}$;
    \item emette $2$ kg di $\mathrm{SO_2}$;
    \item utilizza $50$ litri di petrolio greggio.
\end{itemize}

I due processi non sono indipendenti. Il combustibile prodotto dal
secondo processo viene utilizzato come input dal primo. La LCA deve
tenere conto proprio di queste interdipendenze.


\subsection{La matrice dei processi}

Possiamo raccogliere i vettori dei processi nelle colonne di
un'unica matrice:

\begin{equation}
P=
\left(
p_1\mid p_2
\right)
=
\begin{pmatrix}
-2  & 100\\
10  & 0\\
1   & 10\\
0.1 & 2\\
0   & -50
\end{pmatrix}.
\label{eq:lca-process-matrix}
\end{equation}

Ogni \emph{colonna} della matrice rappresenta quindi un processo,
mentre ogni \emph{riga} rappresenta uno degli item del sistema.

Questa organizzazione è importante dal punto di vista
computazionale: aumentando il numero dei processi, il sistema
produttivo può essere rappresentato aggiungendo nuove colonne;
se compaiono nuovi beni, risorse o emissioni, vengono aggiunte
nuove righe.


\subsection{Tecnosfera e biosfera}

Non tutti i flussi contenuti nella matrice $P$ hanno la stessa
natura.

È utile distinguere tra \emph{tecnosfera} e \emph{biosfera}.

La tecnosfera comprende i beni e i servizi prodotti e utilizzati
all'interno del sistema economico. Nel nostro esempio appartengono
alla tecnosfera:

\begin{equation}
\text{combustibile},
\qquad
\text{elettricità}.
\end{equation}

La biosfera comprende invece gli elementi che rappresentano lo
scambio tra il sistema economico e l'ambiente. Possono essere
risorse prelevate dalla natura oppure sostanze rilasciate
nell'ambiente. Nel nostro esempio appartengono alla biosfera:

\begin{equation}
\mathrm{CO_2},
\qquad
\mathrm{SO_2},
\qquad
\text{petrolio greggio}.
\end{equation}

Possiamo quindi dividere la matrice dei processi in due
sottomatrici:

\begin{equation}
P=
\begin{pmatrix}
A\\
B
\end{pmatrix},
\label{eq:lca-PAB}
\end{equation}

dove $A$ è la \emph{matrice tecnologica} e $B$ è la
\emph{matrice degli interventi ambientali}.

Nel nostro esempio,

\begin{equation}
A=
\begin{pmatrix}
-2 & 100\\
10 & 0
\end{pmatrix},
\end{equation}

mentre

\begin{equation}
B=
\begin{pmatrix}
1   & 10\\
0.1 & 2\\
0   & -50
\end{pmatrix}.
\end{equation}

La matrice $A$ descrive quindi le relazioni tra i processi
all'interno della tecnosfera, mentre $B$ descrive gli scambi tra
i processi produttivi e la biosfera.


\subsection{La domanda finale e l'unità funzionale}

A questo punto possiamo formulare il problema che vogliamo
risolvere.

Supponiamo di essere interessati all'impatto ambientale associato
alla disponibilità di

\begin{equation}
1000\ \text{kWh di elettricità}.
\end{equation}

Questa quantità costituisce, nel nostro esempio, l'\emph{unità
funzionale}: rappresenta il riferimento rispetto al quale vengono
calcolati i flussi e gli impatti ambientali.

Introduciamo il vettore della domanda finale $f$. La sua lunghezza
corrisponde al numero dei flussi economici presenti nella matrice
tecnologica.

Nel nostro caso,

\begin{equation}
f=
\begin{pmatrix}
0\\
1000
\end{pmatrix}.
\label{eq:lca-final-demand}
\end{equation}

Il primo elemento è zero perché non chiediamo al sistema di
fornirci combustibile come prodotto finale. Il combustibile viene
prodotto, ma viene interamente utilizzato come input per la
produzione di elettricità.

Il secondo elemento indica invece che desideriamo ottenere
$1000$ kWh di elettricità.


\subsection{Quante volte devono essere attivati i processi?}

Conosciamo ciò che vogliamo ottenere, ma ogni singolo processo
produce soltanto le quantità indicate nella corrispondente colonna
della matrice $P$.

Dobbiamo quindi determinare quante volte ciascun processo deve
essere attivato affinché il sistema, nel suo complesso, produca
la domanda finale $f$.

Introduciamo a questo scopo il \emph{vettore di scala}

\begin{equation}
s=
\begin{pmatrix}
s_1\\
s_2
\end{pmatrix},
\end{equation}

dove $s_j$ indica il livello di attivazione del processo $j$.

Il vettore di scala deve soddisfare

\begin{equation}
\boxed{
As=f
}.
\label{eq:lca-Asf}
\end{equation}

Ritroviamo quindi esattamente il problema matematico studiato nelle
sezioni precedenti:

\begin{equation}
Ax=b.
\end{equation}

In questo caso, tuttavia, il significato delle variabili è diverso:
la matrice $A$ descrive la tecnologia produttiva, il vettore $f$
rappresenta la domanda finale e il vettore incognito $s$ contiene
i livelli di attività necessari per soddisfarla.

Dal punto di vista puramente algebrico, se $A$ è invertibile,
possiamo scrivere

\begin{equation}
s^*=A^{-1}f.
\label{eq:lca-s-inverse}
\end{equation}

Dal punto di vista computazionale, tuttavia, non è necessario
calcolare esplicitamente $A^{-1}$. Come abbiamo visto nelle
sezioni precedenti, è preferibile considerare direttamente il
problema

\begin{equation}
As=f
\end{equation}

e utilizzare un opportuno metodo numerico per risolverlo.


\subsection{Interpretazione del vettore di scala}

Nel nostro esempio la soluzione è

\begin{equation}
s^*=
\begin{pmatrix}
100\\
2
\end{pmatrix}.
\end{equation}

Questo risultato ha un'interpretazione immediata.

Il primo elemento indica che il processo di produzione
dell'elettricità deve essere attivato $100$ volte. Poiché ogni
attivazione produce $10$ kWh,

\begin{equation}
100\times10=1000\ \text{kWh},
\end{equation}

che corrispondono esattamente alla nostra unità funzionale.

Ma $100$ attivazioni del processo richiedono anche

\begin{equation}
100\times2=200
\end{equation}

litri di combustibile.

Il secondo processo produce $100$ litri di combustibile a ogni
attivazione. Per ottenere i $200$ litri necessari dobbiamo quindi
attivarlo due volte:

\begin{equation}
2\times100=200.
\end{equation}

Il vettore $s^*$ ricostruisce quindi l'intera catena produttiva
necessaria per soddisfare la domanda finale.


\subsection{Dal vettore di scala all'inventario ambientale}

Una volta determinato quante volte devono essere attivati i
processi, possiamo calcolare i corrispondenti scambi con
l'ambiente.

Indichiamo con

\begin{equation}
g=
\begin{pmatrix}
g_1\\
g_2\\
g_3
\end{pmatrix}
\end{equation}

il vettore dei flussi ambientali.

Poiché la matrice $B$ contiene gli interventi ambientali associati
a una singola attivazione dei processi, i flussi complessivi sono

\begin{equation}
\boxed{
g^*=Bs^*
}.
\label{eq:lca-gBs}
\end{equation}

Nel nostro esempio,

\begin{equation}
g^*
=
\begin{pmatrix}
1   & 10\\
0.1 & 2\\
0   & -50
\end{pmatrix}
\begin{pmatrix}
100\\
2
\end{pmatrix}
=
\begin{pmatrix}
120\\
14\\
-100
\end{pmatrix}.
\end{equation}

Il risultato significa che la produzione dei $1000$ kWh di
elettricità richiesti determina complessivamente:

\begin{itemize}
    \item l'emissione di $120$ kg di $\mathrm{CO_2}$;
    \item l'emissione di $14$ kg di $\mathrm{SO_2}$;
    \item il consumo di $100$ litri di petrolio greggio.
\end{itemize}

Il segno negativo dell'ultimo elemento indica che il petrolio
greggio viene prelevato dalla biosfera, mentre i valori positivi
di $\mathrm{CO_2}$ e $\mathrm{SO_2}$ indicano flussi diretti dal
sistema produttivo verso l'ambiente.


\subsection{Dalla domanda finale direttamente ai flussi ambientali}

Poiché

\begin{equation}
s^*=A^{-1}f
\end{equation}

e

\begin{equation}
g^*=Bs^*,
\end{equation}

possiamo combinare le due relazioni:

\begin{equation}
g^*=BA^{-1}f.
\label{eq:lca-direct}
\end{equation}

Definiamo

\begin{equation}
\Lambda=BA^{-1}
\label{eq:lca-intensity}
\end{equation}

come \emph{matrice di intensità}. Otteniamo quindi

\begin{equation}
\boxed{
g^*=\Lambda f
}.
\end{equation}

Questa espressione mette in evidenza la struttura della fase LCI:
la domanda finale determina, attraverso la struttura tecnologica
del sistema, i flussi complessivi che interessano la biosfera.

Anche in questo caso l'espressione con $A^{-1}$ è particolarmente
utile per comprendere la relazione matematica tra domanda finale
e flussi ambientali. In una implementazione numerica, tuttavia,
possiamo ottenere $s^*$ risolvendo direttamente $As=f$ e calcolare
successivamente $g^*=Bs^*$.


\subsection{Il vettore complessivo dei flussi}

Possiamo infine raccogliere la domanda finale e i flussi ambientali
in un unico vettore:

\begin{equation}
q^*=
\begin{pmatrix}
f\\
g^*
\end{pmatrix}.
\end{equation}

Nel nostro esempio,

\begin{equation}
q^*=
\begin{pmatrix}
0\\
1000\\
120\\
14\\
-100
\end{pmatrix}.
\end{equation}

Confrontandolo con il vettore degli item,

\begin{equation}
p_i=
\begin{pmatrix}
\text{litri di combustibile}\\
\text{kWh di elettricità}\\
\text{kg di anidride carbonica}\\
\text{kg di biossido di zolfo}\\
\text{litri di petrolio greggio}
\end{pmatrix},
\end{equation}

possiamo leggere direttamente il risultato dell'inventario.

Alla fine del sistema produttivo disponiamo di $1000$ kWh netti di
elettricità e di nessuna quantità netta di combustibile. Il
combustibile è stato effettivamente prodotto, ma è stato interamente
utilizzato all'interno del sistema per produrre l'elettricità.

Nella biosfera troviamo invece $120$ kg aggiuntivi di
$\mathrm{CO_2}$, $14$ kg aggiuntivi di $\mathrm{SO_2}$ e una
riduzione di $100$ litri di petrolio greggio dovuta al prelievo
della risorsa naturale.


\subsection{Il contributo dei singoli processi}

Il vettore $g^*$ fornisce l'effetto ambientale complessivo del
sistema. Può tuttavia essere utile conoscere anche il contributo
dei singoli processi.

Se indichiamo con $b_j$ la colonna $j$ della matrice $B$, il
contributo ambientale del processo $j$ è

\begin{equation}
b_js_j^*.
\end{equation}

Nel nostro esempio possiamo quindi costruire la matrice

\begin{equation}
\left(
b_1s_1^*
\mid
b_2s_2^*
\right)
=
\begin{pmatrix}
100 & 20\\
10  & 4\\
0   & -100
\end{pmatrix}.
\end{equation}

La stessa informazione può essere rappresentata come

\begin{center}
\begin{tabular}{lrr}
\hline
Flusso ambientale &
Produzione elettricità &
Produzione combustibile\\
\hline
$\mathrm{CO_2}$ & 100 & 20\\
$\mathrm{SO_2}$ & 10 & 4\\
Petrolio greggio & 0 & -100\\
\hline
\end{tabular}
\end{center}

La somma degli elementi di ciascuna riga restituisce il
corrispondente elemento di $g^*$.

Questa rappresentazione permette quindi di distinguere il
contributo dei diversi processi all'inventario ambientale ed è una
forma tipica di output della fase LCI nei software dedicati alla
Life Cycle Assessment.


\subsection{Dall'inventario alla valutazione degli impatti}

La fase LCI ci ha permesso di determinare le quantità fisiche dei
flussi che collegano il sistema produttivo alla biosfera.

Conoscere che la produzione considerata genera $120$ kg di
$\mathrm{CO_2}$ e $14$ kg di $\mathrm{SO_2}$ è già
un'informazione importante, ma non costituisce ancora una misura
degli impatti ambientali.

Sostanze differenti possono infatti contribuire, in misura
differente, a uno stesso problema ambientale. Inoltre, uno stesso
flusso può essere rilevante per differenti categorie di impatto.

La fase di \emph{Life Cycle Impact Assessment} (LCIA) ha quindi
l'obiettivo di trasformare i risultati dell'inventario in indicatori
che rappresentino specifiche categorie di impatto.

Il passaggio può essere sintetizzato come

\begin{equation}
\boxed{
\text{flussi ambientali}
\longrightarrow
\text{categorie di impatto}
\longrightarrow
\text{indicatori di impatto}
}.
\end{equation}


\subsection{Categorie midpoint ed endpoint}

Le categorie di impatto possono essere considerate a differenti
livelli della catena che collega un intervento ambientale al danno
finale.

Le categorie \emph{midpoint} rappresentano effetti intermedi,
relativamente vicini ai flussi dell'inventario. Tra gli esempi
rientrano il cambiamento climatico, l'acidificazione,
l'eutrofizzazione e l'uso delle risorse.

Le categorie \emph{endpoint} cercano invece di rappresentare danni
più vicini alle aree finali che si desidera proteggere.

Tra le principali aree di protezione possiamo considerare:

\begin{itemize}
    \item la salute umana;
    \item l'ambiente naturale;
    \item le risorse naturali.
\end{itemize}

Il passaggio dagli inventari ai midpoint e successivamente agli
endpoint rappresenta quindi una progressiva trasformazione dei
flussi fisici in informazioni sintetiche sugli effetti ambientali.


\subsection{Categorie, indicatori e unità di misura}

Per procedere con la LCIA dobbiamo specificare quale categoria di
impatto vogliamo studiare.

Consideriamo, come primo esempio, il \emph{riscaldamento globale}.

Alla categoria di impatto deve essere associato un indicatore.
Seguendo il semplice esempio considerato qui, assumiamo di
utilizzare come indicatore l'assorbimento della radiazione
infrarossa.

Occorre infine definire un'unità di misura comune attraverso la
quale esprimere il contributo dei differenti flussi ambientali.
Per il riscaldamento globale utilizziamo

\begin{equation}
\text{kg di } \mathrm{CO_2}\text{-equivalente}.
\end{equation}

Abbiamo quindi tre elementi distinti:

\begin{equation}
\boxed{
\text{categoria}
\longrightarrow
\text{indicatore}
\longrightarrow
\text{unità di misura}
}.
\end{equation}


\subsection{Il modello di caratterizzazione}

Per trasformare i flussi dell'inventario nell'indicatore scelto
abbiamo bisogno di un \emph{modello di caratterizzazione}.

Il modello associa a ciascun flusso ambientale un
\emph{fattore di caratterizzazione} (\emph{characterization
factor}, CF).

Un fattore di caratterizzazione misura il contributo di una unità
aggiuntiva di un determinato flusso all'indicatore considerato.

Nell'esempio semplificato adottiamo i seguenti fattori:

\begin{center}
\begin{tabular}{lc}
\hline
Flusso ambientale & Fattore di caratterizzazione\\
\hline
$\mathrm{CO_2}$ & 1\\
$\mathrm{SO_2}$ & 0.1\\
Petrolio greggio & 0\\
\hline
\end{tabular}
\end{center}

Questi valori sono utilizzati esclusivamente per costruire un
esempio numerico semplice. Nelle applicazioni reali, i fattori di
caratterizzazione sono forniti da modelli e metodi di
caratterizzazione sviluppati nella letteratura scientifica.


\subsection{Calcolo di una singola categoria di impatto}

Raccogliamo i fattori di caratterizzazione in un vettore riga:

\begin{equation}
q=
\begin{pmatrix}
1 & 0.1 & 0
\end{pmatrix}.
\end{equation}

Ricordiamo che il vettore dell'inventario è

\begin{equation}
g^*=
\begin{pmatrix}
120\\
14\\
-100
\end{pmatrix}.
\end{equation}

L'indicatore di impatto è ottenuto mediante il prodotto

\begin{equation}
\boxed{
h=qg^*
}.
\label{eq:lca-single-impact}
\end{equation}

Nel nostro esempio,

\begin{align}
h
&=
\begin{pmatrix}
1 & 0.1 & 0
\end{pmatrix}
\begin{pmatrix}
120\\
14\\
-100
\end{pmatrix}\\
&=
120+0.1(14)+0(-100)\\
&=121.4.
\end{align}

L'impatto associato alla categoria considerata è quindi

\begin{equation}
h=121.4
\quad
\text{kg } \mathrm{CO_2}\text{-equivalente}.
\end{equation}

L'operazione di caratterizzazione ha permesso di trasformare flussi
ambientali espressi in unità fisiche differenti in un unico
indicatore espresso in una unità comune.


\subsection{Valutazione simultanea di più categorie di impatto}

Una LCA non si limita normalmente a una sola categoria di impatto.
Supponiamo, ad esempio, di voler considerare simultaneamente:

\begin{equation}
\begin{pmatrix}
\text{acidificazione}\\
\text{riscaldamento globale}\\
\text{deplezione delle risorse}
\end{pmatrix}.
\end{equation}

Per ciascuna categoria dobbiamo individuare:

\begin{itemize}
    \item un indicatore;
    \item una unità di misura;
    \item un modello di caratterizzazione.
\end{itemize}

Supponiamo di utilizzare rispettivamente gli indicatori

\begin{equation}
\begin{pmatrix}
\text{rilascio di }H^+\\
\text{assorbimento della radiazione infrarossa}\\
\text{riduzione della disponibilità della risorsa}
\end{pmatrix}
\end{equation}

e le unità di misura

\begin{equation}
\begin{pmatrix}
\text{kg } \mathrm{SO_2}\text{-equivalente}\\
\text{kg } \mathrm{CO_2}\text{-equivalente}\\
\text{resource depletion units (RDU)}
\end{pmatrix}.
\end{equation}

Per ogni indicatore disponiamo di un insieme di fattori di
caratterizzazione. Possiamo raccoglierli nelle righe di una
matrice $Q$.

Nel nostro esempio,

\begin{equation}
Q=
\begin{pmatrix}
0 & 1   & 0\\
1 & 0.1 & 0\\
0 & 0   & -15
\end{pmatrix}.
\label{eq:lca-Q}
\end{equation}

Le righe rappresentano le categorie di impatto, mentre le colonne
corrispondono ai flussi ambientali:

\begin{equation}
\mathrm{CO_2},
\qquad
\mathrm{SO_2},
\qquad
\text{petrolio greggio}.
\end{equation}

Il vettore degli impatti è quindi

\begin{equation}
\boxed{
h=Qg^*
}.
\label{eq:lca-multiple-impact}
\end{equation}

Nel nostro esempio,

\begin{equation}
h=
\begin{pmatrix}
0 & 1   & 0\\
1 & 0.1 & 0\\
0 & 0   & -15
\end{pmatrix}
\begin{pmatrix}
120\\
14\\
-100
\end{pmatrix}
=
\begin{pmatrix}
14\\
121.4\\
1500
\end{pmatrix}.
\end{equation}

Otteniamo quindi simultaneamente:

\begin{itemize}
    \item $14$ kg $\mathrm{SO_2}$-equivalente per
          l'acidificazione;
    \item $121.4$ kg $\mathrm{CO_2}$-equivalente per il
          riscaldamento globale;
    \item $1500$ RDU per la deplezione delle risorse.
\end{itemize}


\subsection{La struttura computazionale completa della LCA}

Possiamo ora riunire i diversi passaggi.

Partiamo dalla domanda finale

\begin{equation}
f.
\end{equation}

La struttura tecnologica del sistema determina i livelli di
attività necessari:

\begin{equation}
As=f.
\end{equation}

Risolvendo il sistema otteniamo

\begin{equation}
s^*=A^{-1}f.
\end{equation}

Il vettore di scala permette di determinare l'inventario
ambientale:

\begin{equation}
g^*=Bs^*.
\end{equation}

Sostituendo la relazione precedente possiamo scrivere

\begin{equation}
g^*=BA^{-1}f.
\end{equation}

Infine, la matrice dei fattori di caratterizzazione trasforma i
flussi ambientali negli indicatori di impatto:

\begin{equation}
h=Qg^*.
\end{equation}

L'intera struttura può quindi essere sintetizzata come

\begin{equation}
\boxed{
\begin{array}{ccccc}
f
&\longrightarrow&
s^*
&\longrightarrow&
g^*
\longrightarrow h\\[2mm]
\text{domanda finale}
&&
\text{livelli di attività}
&&
\text{inventario}
\quad
\text{impatti}
\end{array}}
\end{equation}

oppure, evidenziando le operazioni matematiche,

\begin{equation}
\boxed{
As=f,
\qquad
g^*=Bs^*,
\qquad
h=Qg^*
}.
\label{eq:lca-computational-structure}
\end{equation}

Questa rappresentazione permette di comprendere perché la LCA
costituisca un'interessante applicazione degli strumenti di algebra
lineare studiati in precedenza. Il problema centrale della fase
di inventario consiste ancora una volta nella soluzione di un
sistema di equazioni lineari.

Cambia il contesto applicativo, ma non la struttura matematica:

\begin{equation}
\boxed{
\begin{array}{ccc}
\text{modello economico} &:& Ax=b,\\[1mm]
\text{Life Cycle Inventory} &:& As=f.
\end{array}}
\end{equation}

Una volta determinata la soluzione del sistema lineare, ulteriori
prodotti matrice-vettore permettono di passare dai livelli di
attività ai flussi ambientali e, successivamente, dai flussi
ambientali agli indicatori di impatto.

La sequenza

\begin{equation}
\boxed{
\text{domanda finale}
\longrightarrow
\text{inventario}
\longrightarrow
\text{impatti}
}
\end{equation}

riassume quindi la struttura computazionale essenziale della
Life Cycle Assessment considerata in questa sezione.

\chapter{La struttura a termine dei tassi di interesse}

\section{Richiami teorici}
\label{sec:term-structure}

La terza applicazione dei metodi numerici per la soluzione dei
sistemi lineari riguarda la finanza e, in particolare, la
costruzione della \emph{struttura a termine dei tassi di interesse}.

Il problema che vogliamo affrontare può essere formulato nel modo
seguente: osservando sul mercato i prezzi di un insieme di
obbligazioni e conoscendo i flussi di cassa che esse genereranno
in futuro, possiamo ricavare i tassi di interesse che il mercato
associa alle diverse scadenze?

Come vedremo, anche questo problema può essere ricondotto alla
soluzione di un sistema di equazioni lineari. In questo modo,
lo stesso problema numerico

\begin{equation}
    Ax=b
\end{equation}

che abbiamo incontrato nelle applicazioni economiche e nella
Life Cycle Assessment ricompare in un contesto finanziario.


\subsection{Valore del denaro nel tempo e fattori di sconto}

Un principio fondamentale della finanza è che il valore di una
somma di denaro dipende dal momento nel quale essa diventa
disponibile. In generale, un euro disponibile oggi vale più di un
euro che verrà ricevuto in futuro.

Supponiamo di ricevere con certezza un ammontare $FV$ al tempo
$T$. Se il tasso di interesse relativo a quella scadenza è $r_T$,
il valore attuale del pagamento futuro è

\begin{equation}
    PV=
    \frac{FV}{(1+r_T)^T}.
    \label{eq:present-value-term}
\end{equation}

Possiamo separare il pagamento futuro dal termine che permette
di riportarlo al suo valore attuale definendo il
\emph{fattore di sconto}

\begin{equation}
    d_T=
    \frac{1}{(1+r_T)^T}.
    \label{eq:discount-factor}
\end{equation}

L'Equazione~\eqref{eq:present-value-term} può quindi essere
riscritta semplicemente come

\begin{equation}
    PV=d_TFV.
\end{equation}

Il fattore di sconto $d_T$ ha un'interpretazione finanziaria
immediata: rappresenta il valore oggi di una unità monetaria che
verrà ricevuta con certezza al tempo $T$.

Ad esempio, se

\begin{equation}
    d_T=0.90,
\end{equation}

una unità monetaria disponibile al tempo $T$ ha oggi un valore
pari a $0.90$.

Tasso di interesse e fattore di sconto rappresentano quindi due
modi differenti di esprimere la stessa relazione tra valori
monetari disponibili in momenti diversi.


\subsection{La struttura a termine dei tassi di interesse}

Finora abbiamo indicato il tasso di interesse con $r_T$,
specificando esplicitamente la scadenza $T$. Questa indicazione
è importante perché sui mercati finanziari non si utilizza
necessariamente lo stesso tasso per valutare pagamenti che si
verificano in date differenti.

In generale,

\begin{equation}
    r_1\neq r_2\neq\cdots\neq r_T.
\end{equation}

Un pagamento che avverrà tra un anno può quindi essere scontato
utilizzando un tasso differente da quello utilizzato per un
pagamento che avverrà tra due, cinque o dieci anni.

La relazione che associa a ogni scadenza il corrispondente tasso
di interesse,

\begin{equation}
    T\longmapsto r_T,
\end{equation}

prende il nome di \emph{struttura a termine dei tassi di
interesse}.

La sua rappresentazione grafica è comunemente chiamata
\emph{curva dei rendimenti} o \emph{yield curve}.

I tassi $r_T$ utilizzati per scontare singoli pagamenti alle
diverse scadenze sono chiamati \emph{tassi spot} o
\emph{tassi zero-coupon}.


\subsection{Tassi spot e rendimento a scadenza}

È importante distinguere il tasso spot dal
\emph{rendimento a scadenza} di un'obbligazione
(\emph{yield to maturity}, YTM).

Il tasso spot $r_T$ è il tasso appropriato per scontare un
singolo pagamento che si verifica al tempo $T$:

\begin{equation}
    PV=
    \frac{FV}{(1+r_T)^T}.
\end{equation}

Consideriamo invece un'obbligazione che genera diversi flussi di
cassa

\begin{equation}
    CF_1,CF_2,\ldots,CF_T.
\end{equation}

Il rendimento a scadenza $y$ è definito come l'unico tasso che,
applicato a tutti i flussi di cassa dell'obbligazione, rende il
loro valore attuale uguale al prezzo di mercato $P$:

\begin{equation}
    P=
    \sum_{t=1}^{T}
    \frac{CF_t}{(1+y)^t}.
    \label{eq:ytm}
\end{equation}

La differenza è quindi sostanziale.

Nel calcolo dello YTM utilizziamo lo stesso tasso $y$ per tutti
i flussi di cassa dell'obbligazione. Nella struttura a termine,
invece, ogni scadenza può essere associata a un proprio tasso
spot:

\begin{equation}
    r_1,r_2,\ldots,r_T.
\end{equation}

Il rendimento a scadenza sintetizza quindi in un unico numero il
rendimento associato all'intera obbligazione, mentre i tassi spot
descrivono separatamente il valore temporale del denaro alle
diverse scadenze.

Poiché il nostro obiettivo è costruire la struttura a termine,
siamo interessati ai tassi spot $r_t$ e ai corrispondenti fattori
di sconto $d_t$, e non direttamente allo YTM dei singoli titoli.


\subsection{Valutazione di un'obbligazione mediante fattori di sconto}

Consideriamo ora un'obbligazione che genera i flussi di cassa

\begin{equation}
    CF_1,CF_2,\ldots,CF_T.
\end{equation}

Poiché i pagamenti avvengono in date differenti, ciascun flusso
deve essere valutato utilizzando il fattore di sconto
corrispondente alla propria scadenza.

Il prezzo dell'obbligazione è quindi

\begin{equation}
    P=
    CF_1d_1+
    CF_2d_2+
    \cdots+
    CF_Td_T,
\end{equation}

oppure, in forma compatta,

\begin{equation}
    \boxed{
    P=
    \sum_{t=1}^{T}CF_td_t
    }.
    \label{eq:bond-discount-factors}
\end{equation}

Questa relazione è fondamentale per il problema che vogliamo
studiare.

Se conosciamo i fattori di sconto, possiamo infatti utilizzarli
per calcolare il valore attuale dei flussi di cassa e quindi il
prezzo dell'obbligazione.


\subsection{Un semplice esempio}

Consideriamo un'obbligazione con scadenza a tre anni, valore
nominale pari a $100$ e cedola annuale pari a $5$.

I flussi di cassa sono

\begin{equation}
    CF_1=5,
    \qquad
    CF_2=5,
    \qquad
    CF_3=105.
\end{equation}

Nei primi due anni l'investitore riceve soltanto la cedola. Alla
scadenza riceve invece l'ultima cedola, pari a $5$, insieme al
rimborso del valore nominale, pari a $100$.

Il prezzo dell'obbligazione è pertanto

\begin{equation}
    P=5d_1+5d_2+105d_3.
    \label{eq:bond-three-years}
\end{equation}

L'esempio mette in evidenza un aspetto importante: il prezzo di
un'obbligazione con cedole dipende generalmente da più fattori di
sconto, perché i suoi flussi di cassa si verificano in date
differenti.


\subsection{Il problema inverso: dai prezzi ai fattori di sconto}

L'Equazione~\eqref{eq:bond-discount-factors} può essere utilizzata
in due direzioni.

Se conosciamo i fattori di sconto

\begin{equation}
    d_1,d_2,\ldots,d_T,
\end{equation}

e conosciamo i flussi di cassa dell'obbligazione, possiamo
calcolarne il prezzo.

Il problema che incontriamo nella realtà è però spesso quello
opposto.

I prezzi delle obbligazioni sono osservabili sui mercati
finanziari. Anche i flussi di cassa contrattualmente previsti
dalle obbligazioni sono noti. I fattori di sconto, invece, non
sono direttamente osservabili.

Possiamo allora chiederci se sia possibile utilizzare le
informazioni osservabili per ricavare quelle non osservabili.

Supponiamo di osservare $N$ obbligazioni e di considerare $N$
diverse scadenze. Per ciascuna obbligazione $i$ possiamo scrivere

\begin{equation}
    P_i=
    \sum_{t=1}^{N}CF_{it}d_t,
    \qquad
    i=1,\ldots,N,
    \label{eq:bond-general-system}
\end{equation}

dove $CF_{it}$ rappresenta il flusso di cassa generato
dall'obbligazione $i$ alla scadenza $t$.

Ogni obbligazione fornisce quindi un'equazione. Osservando più
obbligazioni otteniamo un sistema di equazioni nelle incognite

\begin{equation}
    d_1,d_2,\ldots,d_N.
\end{equation}

Il problema finanziario può quindi essere formulato nel modo
seguente:

\begin{quote}
è possibile ricavare la struttura a termine dei tassi di interesse
utilizzando i prezzi delle obbligazioni osservati sul mercato?
\end{quote}


\subsection{Un esempio con tre obbligazioni}

Consideriamo tre obbligazioni, indicate con $A$, $B$ e $C$, con
i seguenti flussi di cassa e prezzi:

\begin{center}
\begin{tabular}{lrrrr}
\hline
Obbligazione & $t=1$ & $t=2$ & $t=3$ & Prezzo\\
\hline
A & 105 & 0   & 0   & 100\\
B & 4   & 104 & 0   & 99\\
C & 5   & 5   & 105 & 98\\
\hline
\end{tabular}
\end{center}

Per l'obbligazione $A$ abbiamo

\begin{equation}
    105d_1=100.
    \label{eq:bond-A}
\end{equation}

Per l'obbligazione $B$:

\begin{equation}
    4d_1+104d_2=99.
    \label{eq:bond-B}
\end{equation}

Infine, per l'obbligazione $C$:

\begin{equation}
    5d_1+5d_2+105d_3=98.
    \label{eq:bond-C}
\end{equation}

Le tre incognite sono quindi i fattori di sconto

\begin{equation}
    d=
    \begin{pmatrix}
        d_1\\
        d_2\\
        d_3
    \end{pmatrix}.
\end{equation}


\subsection{Dal problema finanziario al sistema lineare}

Le tre equazioni precedenti possono essere raccolte in forma
matriciale:

\begin{equation}
\begin{pmatrix}
105 & 0   & 0\\
4   & 104 & 0\\
5   & 5   & 105
\end{pmatrix}
\begin{pmatrix}
d_1\\
d_2\\
d_3
\end{pmatrix}
=
\begin{pmatrix}
100\\
99\\
98
\end{pmatrix}.
\label{eq:bond-matrix-example}
\end{equation}

Definiamo

\begin{equation}
C=
\begin{pmatrix}
105 & 0   & 0\\
4   & 104 & 0\\
5   & 5   & 105
\end{pmatrix}
\end{equation}

come matrice dei flussi di cassa,

\begin{equation}
d=
\begin{pmatrix}
d_1\\
d_2\\
d_3
\end{pmatrix}
\end{equation}

come vettore dei fattori di sconto e

\begin{equation}
P=
\begin{pmatrix}
100\\
99\\
98
\end{pmatrix}
\end{equation}

come vettore dei prezzi osservati.

Il problema può quindi essere scritto nella forma compatta

\begin{equation}
    \boxed{
    Cd=P
    }.
    \label{eq:CdP}
\end{equation}

Abbiamo così trasformato il problema finanziario in un sistema di
equazioni lineari.

La corrispondenza con la notazione generale utilizzata nelle
sezioni precedenti è immediata:

\begin{equation}
\boxed{
\begin{array}{ccc}
Ax=b
&\longleftrightarrow&
Cd=P.
\end{array}}
\end{equation}

La matrice dei coefficienti $A$ è ora la matrice dei flussi di
cassa $C$; il vettore delle incognite $x$ è il vettore dei fattori
di sconto $d$; il vettore dei termini noti $b$ è costituito dai
prezzi osservati $P$.

Questo passaggio è particolarmente importante perché mostra che
una tecnica numerica generale può essere utilizzata per risolvere
problemi appartenenti ad ambiti applicativi molto differenti.


\subsection{Bootstrapping della struttura a termine}

Nell'esempio appena considerato la matrice dei flussi di cassa è
triangolare inferiore:

\begin{equation}
C=
\begin{pmatrix}
105 & 0   & 0\\
4   & 104 & 0\\
5   & 5   & 105
\end{pmatrix}.
\end{equation}

Questa particolare struttura permette di determinare i fattori di
sconto uno dopo l'altro.

Dalla prima equazione,

\begin{equation}
    105d_1=100,
\end{equation}

otteniamo immediatamente

\begin{equation}
    d_1=\frac{100}{105}.
    \label{eq:d1-bootstrap}
\end{equation}

Una volta noto $d_1$, possiamo utilizzare la seconda equazione:

\begin{equation}
    4d_1+104d_2=99.
\end{equation}

Isolando $d_2$ otteniamo

\begin{equation}
    d_2=
    \frac{99-4d_1}{104}.
    \label{eq:d2-bootstrap}
\end{equation}

A questo punto conosciamo $d_1$ e $d_2$ e possiamo utilizzare la
terza equazione:

\begin{equation}
    5d_1+5d_2+105d_3=98,
\end{equation}

da cui

\begin{equation}
    d_3=
    \frac{98-5d_1-5d_2}{105}.
    \label{eq:d3-bootstrap}
\end{equation}

Abbiamo quindi ricavato progressivamente

\begin{equation}
    d_1
    \longrightarrow
    d_2
    \longrightarrow
    d_3.
\end{equation}

Questa procedura sequenziale prende il nome di
\emph{bootstrapping}.

Nel nostro esempio il bootstrapping è particolarmente semplice
perché la matrice $C$ è triangolare. Dal punto di vista numerico,
stiamo quindi utilizzando la stessa logica della sostituzione in
avanti incontrata nello studio dei sistemi triangolari.

Più in generale, la struttura dei flussi di cassa potrebbe non
permettere una soluzione sequenziale così immediata. Il problema
rimane tuttavia rappresentabile come

\begin{equation}
    Cd=P
\end{equation}

e può essere affrontato utilizzando le tecniche numeriche per la
soluzione dei sistemi lineari introdotte nelle sezioni precedenti.


\subsection{Dai fattori di sconto ai tassi spot}

La soluzione del sistema ci fornisce i fattori di sconto

\begin{equation}
    d_1,d_2,\ldots,d_T.
\end{equation}

Il nostro obiettivo finale è però costruire la struttura a termine
dei tassi di interesse. Dobbiamo quindi trasformare ciascun fattore
di sconto nel corrispondente tasso spot.

Ricordiamo la relazione

\begin{equation}
    d_t=
    \frac{1}{(1+r_t)^t}.
\end{equation}

Elevando entrambi i membri alla potenza $-1/t$ otteniamo

\begin{equation}
    1+r_t=d_t^{-1/t},
\end{equation}

e quindi

\begin{equation}
    \boxed{
    r_t=d_t^{-1/t}-1
    }.
    \label{eq:discount-to-spot}
\end{equation}

Una volta calcolato un fattore di sconto per ciascuna scadenza,
possiamo pertanto ottenere i corrispondenti tassi spot:

\begin{equation}
\begin{array}{ccccccc}
d_1 & d_2 & d_3 & \cdots & d_T\\
\downarrow & \downarrow & \downarrow && \downarrow\\
r_1 & r_2 & r_3 & \cdots & r_T.
\end{array}
\end{equation}

Associando ciascun tasso $r_t$ alla propria scadenza $t$ otteniamo
la struttura a termine

\begin{equation}
    t\longmapsto r_t,
\end{equation}

che può essere rappresentata graficamente attraverso la curva dei
rendimenti.


\subsection{La struttura computazionale del problema}

Possiamo infine riassumere il procedimento seguito.

I dati inizialmente disponibili sono i prezzi delle obbligazioni
e i rispettivi flussi di cassa:

\begin{equation}
    C,\qquad P.
\end{equation}

Da queste informazioni costruiamo il sistema

\begin{equation}
    Cd=P.
\end{equation}

La soluzione del sistema fornisce i fattori di sconto:

\begin{equation}
    C,P
    \longrightarrow
    d.
\end{equation}

Successivamente trasformiamo i fattori di sconto nei
corrispondenti tassi spot:

\begin{equation}
    d_t
    \longrightarrow
    r_t=d_t^{-1/t}-1.
\end{equation}

L'intero procedimento può quindi essere rappresentato come

\begin{equation}
\boxed{
\begin{array}{c}
\text{prezzi e flussi di cassa}\\
\downarrow\\
Cd=P\\
\downarrow\\
\text{fattori di sconto}\\
\downarrow\\
r_t=d_t^{-1/t}-1\\
\downarrow\\
\text{tassi spot}\\
\downarrow\\
\text{curva dei rendimenti}
\end{array}}
\end{equation}

Il problema finanziario della costruzione della struttura a termine
si riconduce quindi, nella sua parte centrale, allo stesso problema
numerico

\begin{equation}
    Ax=b
\end{equation}

incontrato nelle precedenti applicazioni del corso.

\section{Richiami: titoli di Stato italiani}
\label{subsec:titoli-stato-italiani}

Finora abbiamo studiato il problema della costruzione della struttura
a termine utilizzando obbligazioni teoriche, caratterizzate da flussi
di cassa particolarmente semplici.

L'obiettivo è ora applicare la stessa procedura a dati osservati sui
mercati finanziari. Prima di procedere con l'applicazione numerica è
quindi utile richiamare brevemente le principali caratteristiche dei
titoli di Stato italiani.

I titoli di Stato sono strumenti finanziari emessi dal Ministero
dell'Economia e delle Finanze per finanziare il debito pubblico.
Dal punto di vista dell'investitore, l'acquisto di un titolo comporta
il pagamento di un prezzo al momento dell'acquisto e il diritto a
ricevere uno o più flussi di cassa futuri.

La struttura di questi flussi non è però uguale per tutti i titoli.
Alcuni strumenti non prevedono il pagamento di cedole, mentre altri
generano pagamenti periodici. Esistono inoltre titoli nei quali le
cedole o il capitale rimborsato dipendono dall'andamento di un
determinato indice.

Queste differenze sono importanti per la nostra applicazione.
Per costruire la matrice dei flussi di cassa $C$ dobbiamo infatti
conoscere, per ciascun titolo, non soltanto il prezzo osservato, ma
anche l'ammontare e la data di ciascun pagamento futuro.


\subsection{Buoni Ordinari del Tesoro}

I \emph{Buoni Ordinari del Tesoro} (BOT) sono titoli a breve termine
che non prevedono il pagamento di cedole.

Il rendimento dell'investitore deriva dalla differenza tra il prezzo
pagato per acquistare il titolo e il valore rimborsato alla scadenza.

In forma semplificata, la struttura dei flussi può essere
rappresentata come

\begin{equation}
\boxed{
-P_0
\qquad\longrightarrow\qquad
F_T
}
\end{equation}

dove $P_0$ rappresenta il prezzo pagato oggi e $F_T$ il valore
rimborsato alla scadenza.

Dal punto di vista della costruzione della struttura a termine,
questa caratteristica rende i BOT particolarmente semplici.

Se un titolo genera un unico pagamento $F_T$ alla scadenza $T$,
il suo prezzo soddisfa infatti

\begin{equation}
    P_0=F_Td_T.
\end{equation}

Il fattore di sconto relativo a quella scadenza può quindi essere
ricavato direttamente:

\begin{equation}
    \boxed{
    d_T=\frac{P_0}{F_T}
    }.
\end{equation}

I titoli privi di cedole costituiscono pertanto il caso più semplice
per osservare la relazione tra prezzo, fattore di sconto e tasso
spot.


\subsection{Buoni del Tesoro Poliennali}

I \emph{Buoni del Tesoro Poliennali} (BTP) sono titoli caratterizzati
da una scadenza più lunga e dal pagamento periodico di cedole.

Indicando con $c_t$ la cedola pagata alla data $t$ e con $F$ il
valore nominale rimborsato alla scadenza, la sequenza dei flussi di
cassa può essere rappresentata schematicamente come

\begin{equation}
\boxed{
-P_0
\longrightarrow
c_1
\longrightarrow
c_2
\longrightarrow
\cdots
\longrightarrow
c_T+F
}.
\end{equation}

Il prezzo del titolo dipende quindi dal valore attuale di tutti i
pagamenti futuri:

\begin{equation}
    P_0
    =
    c_1d_1+c_2d_2+\cdots+(c_T+F)d_T.
    \label{eq:btp-price}
\end{equation}

Questa struttura è precisamente quella che abbiamo incontrato
nell'esempio teorico precedente.

A differenza di un BOT, quindi, un singolo BTP non permette
generalmente di determinare direttamente un unico fattore di sconto:
il suo prezzo dipende contemporaneamente dai fattori di sconto
associati a tutte le date nelle quali il titolo genera un pagamento.

È proprio questa caratteristica a rendere naturale la costruzione
del sistema

\begin{equation}
    Cd=P.
\end{equation}


\subsection{Certificati di Credito del Tesoro}

I \emph{Certificati di Credito del Tesoro} (CCT) appartengono alla
categoria dei titoli a tasso variabile.

A differenza di un titolo a cedola fissa, l'ammontare delle cedole
future non è necessariamente noto al momento dell'emissione, ma
dipende dal parametro al quale il titolo è indicizzato.

La struttura generale dei flussi rimane

\begin{equation}
-P_0
\longrightarrow
c_1
\longrightarrow
c_2
\longrightarrow
\cdots
\longrightarrow
c_T+F,
\end{equation}

ma le cedole

\begin{equation}
    c_1,c_2,\ldots,c_T
\end{equation}

non sono tutte fissate anticipatamente nello stesso modo di un
titolo a tasso fisso.

Questa differenza è rilevante per la costruzione della matrice dei
flussi di cassa: per utilizzare un titolo nella procedura descritta
precedentemente dobbiamo essere in grado di determinare i pagamenti
associati alle diverse date.


\subsection{Titoli indicizzati all'inflazione}

Un'ulteriore categoria è costituita dai titoli il cui rendimento è
collegato all'andamento dell'inflazione.

In questi strumenti l'indicizzazione ha l'obiettivo di collegare
almeno una parte dei pagamenti ricevuti dall'investitore
all'evoluzione del livello dei prezzi.

Dal punto di vista della nostra applicazione, la caratteristica
importante è che i flussi di cassa futuri non possono essere
trattati nello stesso modo dei flussi di un'obbligazione nominale
a cedola fissa.

Per costruire una struttura a termine nominale mediante il sistema

\begin{equation}
    Cd=P,
\end{equation}

è infatti particolarmente conveniente lavorare con strumenti per
i quali i flussi di cassa nominali futuri siano determinabili.


\subsection{Perché la tipologia del titolo è importante}

Il breve richiamo alle diverse tipologie di titoli di Stato permette
di comprendere un punto importante per la successiva applicazione.

Il prezzo di mercato, da solo, non è sufficiente per costruire la
struttura a termine. Per ogni titolo dobbiamo conoscere l'intera
sequenza dei flussi di cassa:

\begin{equation}
    CF_{i1},CF_{i2},\ldots,CF_{iT}.
\end{equation}

Questi valori costituiscono le righe della matrice

\begin{equation}
C=
\begin{pmatrix}
CF_{11} & CF_{12} & \cdots & CF_{1T}\\
CF_{21} & CF_{22} & \cdots & CF_{2T}\\
\vdots  & \vdots  &        & \vdots\\
CF_{N1} & CF_{N2} & \cdots & CF_{NT}
\end{pmatrix}.
\end{equation}

Il vettore

\begin{equation}
P=
\begin{pmatrix}
P_1\\
P_2\\
\vdots\\
P_N
\end{pmatrix}
\end{equation}

contiene invece i prezzi osservati sul mercato.

Il problema finanziario diventa ancora una volta

\begin{equation}
    Cd=P.
\end{equation}

Nell'esempio teorico precedente avevamo costruito artificialmente
titoli con date dei pagamenti perfettamente ordinate e una matrice
dei flussi di cassa triangolare.

Quando utilizziamo titoli realmente negoziati sul mercato, la
situazione è inevitabilmente più articolata. I titoli possono avere
date di scadenza differenti, cedole differenti e calendari dei
pagamenti che non coincidono perfettamente.

Prima di poter applicare gli strumenti numerici studiati nel corso
sarà quindi necessario organizzare opportunamente queste
informazioni.

Questo passaggio ci porterà dall'esempio teorico alla costruzione
della struttura a termine utilizzando dati finanziari reali.

\section{Interpolazione numerica}
\label{sec:interpolazione}

Nella sezione precedente abbiamo visto che i prezzi dei titoli di Stato
possono essere utilizzati per ricavare informazioni sulla struttura a
termine dei fattori di sconto. In particolare, i BOT sono strumenti
particolarmente semplici perché generano un unico pagamento alla
scadenza. Se il valore nominale è pari a 100, dal prezzo di un BOT
possiamo ricavare direttamente il fattore di sconto:

\begin{equation}
    d_T = \frac{P_T}{100}.
\end{equation}

Questo risultato sembra fornire una soluzione immediata al nostro
problema. Esiste tuttavia una difficoltà importante quando passiamo
dall'esempio teorico ai dati realmente osservati sul mercato.

I BOT disponibili non hanno una scadenza per ogni possibile data
futura. Essi permettono quindi di conoscere direttamente la funzione
di sconto soltanto in corrispondenza di un insieme discreto di
scadenze.

Supponiamo, ad esempio, di conoscere i fattori di sconto

\begin{equation}
    d(T_1), d(T_2), \ldots, d(T_n)
\end{equation}

corrispondenti alle scadenze

\begin{equation}
    T_1,T_2,\ldots,T_n.
\end{equation}

Potremmo però avere la necessità di valutare un pagamento che avviene
a una data $T^*$ che non coincide con nessuna di queste scadenze.

Il problema diventa allora

\begin{equation}
\boxed{
\begin{array}{c}
\text{conosciamo } d(T_1),d(T_2),\ldots,d(T_n)\\[2mm]
\downarrow\\[2mm]
\text{vogliamo determinare } d(T^*)
\end{array}}
\end{equation}

per una scadenza $T^*$ non direttamente osservata.

Questo è un problema di \emph{interpolazione}.

L'interpolazione consiste nel costruire una funzione che passa
attraverso un insieme di punti noti e nell'utilizzare tale funzione
per stimare il valore corrispondente a punti intermedi.

Nel nostro caso i punti osservati sono costituiti dalle coppie

\begin{equation}
    (T_i,d_i),
\end{equation}

dove $T_i$ rappresenta la maturity e $d_i$ il corrispondente fattore
di sconto.


\subsection{Interpolazione lineare}

Consideriamo inizialmente il caso più semplice. Supponiamo di conoscere
due punti della funzione di sconto:

\begin{equation}
    (T_1,d_1)
    \qquad\text{e}\qquad
    (T_2,d_2),
\end{equation}

e di voler stimare il fattore di sconto relativo a una maturity
intermedia $T^*$, con

\begin{equation}
    T_1<T^*<T_2.
\end{equation}

L'ipotesi più semplice consiste nell'assumere che, nell'intervallo
compreso tra $T_1$ e $T_2$, la funzione di sconto possa essere
approssimata mediante una retta:

\begin{equation}
    d(T)=a+bT.
    \label{eq:linear-interpolation}
\end{equation}

I coefficienti $a$ e $b$ non sono noti. Tuttavia conosciamo due punti
attraverso i quali la retta deve necessariamente passare.

Imponendo

\begin{equation}
    d(T_1)=d_1
\end{equation}

otteniamo

\begin{equation}
    a+bT_1=d_1.
\end{equation}

Analogamente, imponendo

\begin{equation}
    d(T_2)=d_2
\end{equation}

otteniamo

\begin{equation}
    a+bT_2=d_2.
\end{equation}

Abbiamo quindi due equazioni nelle due incognite $a$ e $b$:

\begin{equation}
\begin{cases}
a+bT_1=d_1,\\
a+bT_2=d_2.
\end{cases}
\end{equation}

Anche il problema dell'interpolazione può quindi essere ricondotto
alla soluzione di un sistema lineare.


\subsection{Interpolazione lineare come sistema di equazioni}

Le due equazioni precedenti possono essere scritte in forma
matriciale:

\begin{equation}
\begin{pmatrix}
1 & T_1\\
1 & T_2
\end{pmatrix}
\begin{pmatrix}
a\\
b
\end{pmatrix}
=
\begin{pmatrix}
d_1\\
d_2
\end{pmatrix}.
\label{eq:linear-interpolation-system}
\end{equation}

Ritroviamo ancora una volta la struttura generale

\begin{equation}
    Ax=b.
\end{equation}

In questo caso

\begin{equation}
A=
\begin{pmatrix}
1&T_1\\
1&T_2
\end{pmatrix},
\qquad
x=
\begin{pmatrix}
a\\
b
\end{pmatrix},
\qquad
b=
\begin{pmatrix}
d_1\\
d_2
\end{pmatrix}.
\end{equation}

Risolvendo il sistema determiniamo i coefficienti $a$ e $b$ della
retta. A quel punto possiamo calcolare il fattore di sconto
corrispondente a qualsiasi maturity $T^*$ compresa nell'intervallo:

\begin{equation}
    d(T^*)=a+bT^*.
\end{equation}

È importante osservare il ruolo svolto dal sistema lineare.
Le incognite non sono direttamente i fattori di sconto che vogliamo
interpolare, ma i coefficienti della funzione utilizzata per
rappresentare la relazione tra maturity e fattore di sconto.

Una volta determinata questa funzione, essa può essere valutata in
qualunque punto dell'intervallo considerato.


\subsection{Un esempio con dati finanziari reali}

Il problema appena descritto emerge direttamente nei dati relativi
ai titoli di Stato.

Consideriamo il titolo

\begin{center}
\texttt{BTP 01/11/2026 7,25\%}.
\end{center}

Per valutare il pagamento che avviene alla scadenza del titolo abbiamo
bisogno del fattore di sconto relativo al 1 novembre 2026.

Supponiamo però che nel dataset non sia presente un BOT con questa
esatta scadenza. I due BOT con scadenze più vicine sono:

\begin{center}
\begin{tabular}{lll}
\hline
Titolo & Scadenza & Fattore di sconto\\
\hline
BOT 14/10/2026 ZC & 14/10/2026 & $d_1$\\
BOT 13/11/2026 ZC & 13/11/2026 & $d_2$\\
\hline
\end{tabular}
\end{center}

La data del pagamento del BTP si trova tra le due date:

\begin{equation}
14/10/2026
<
01/11/2026
<
13/11/2026.
\end{equation}

Possiamo quindi utilizzare i fattori di sconto dei due BOT per
costruire una retta di interpolazione e valutare tale retta in
corrispondenza del 1 novembre 2026.

In termini generali, il procedimento è

\begin{equation}
\boxed{
(T_1,d_1),(T_2,d_2)
\longrightarrow
\begin{pmatrix}
1&T_1\\
1&T_2
\end{pmatrix}
\begin{pmatrix}
a\\
b
\end{pmatrix}
=
\begin{pmatrix}
d_1\\
d_2
\end{pmatrix}
\longrightarrow
d(T^*)
}.
\end{equation}

L'esempio mostra perché l'interpolazione è necessaria nella pratica:
le date dei pagamenti dei BTP non coincidono necessariamente con le
scadenze dei BOT dai quali possiamo osservare direttamente i fattori
di sconto.


\subsection{Dall'interpolazione lineare all'interpolazione polinomiale}

L'interpolazione lineare utilizza due punti osservati e costruisce una
retta che li collega:

\begin{equation}
    d(T)=a_0+a_1T.
\end{equation}

Possiamo però generalizzare questa idea.

Supponiamo, ad esempio, di conoscere tre punti:

\begin{equation}
    (T_1,d_1),\qquad
    (T_2,d_2),\qquad
    (T_3,d_3).
\end{equation}

Invece di utilizzare una retta possiamo cercare un polinomio di
secondo grado:

\begin{equation}
    d(T)=a_0+a_1T+a_2T^2.
\end{equation}

Imponendo che il polinomio passi attraverso i tre punti otteniamo

\begin{equation}
\begin{cases}
a_0+a_1T_1+a_2T_1^2=d_1,\\
a_0+a_1T_2+a_2T_2^2=d_2,\\
a_0+a_1T_3+a_2T_3^2=d_3.
\end{cases}
\end{equation}

Abbiamo quindi tre equazioni nelle tre incognite

\begin{equation}
    a_0,\qquad a_1,\qquad a_2.
\end{equation}

La stessa idea può essere estesa a un numero arbitrario di punti.


\subsection{Interpolazione polinomiale}

Consideriamo $n$ BOT dai quali abbiamo ricavato $n$ punti della
funzione di sconto:

\begin{equation}
    (T_1,d_1),(T_2,d_2),\ldots,(T_n,d_n).
\end{equation}

Cerchiamo ora un polinomio di grado $n-1$:

\begin{equation}
    d(T)
    =
    a_0+a_1T+a_2T^2+\cdots+a_{n-1}T^{n-1}.
    \label{eq:polynomial-interpolation}
\end{equation}

Il polinomio contiene esattamente $n$ coefficienti incogniti:

\begin{equation}
    a_0,a_1,\ldots,a_{n-1}.
\end{equation}

Per determinarli imponiamo che il polinomio passi attraverso tutti i
punti osservati:

\begin{equation}
    d(T_i)=d_i,
    \qquad i=1,\ldots,n.
\end{equation}

Otteniamo quindi il sistema

\begin{equation}
\begin{cases}
a_0+a_1T_1+a_2T_1^2+\cdots+a_{n-1}T_1^{n-1}=d_1,\\
a_0+a_1T_2+a_2T_2^2+\cdots+a_{n-1}T_2^{n-1}=d_2,\\
\hspace{2cm}\vdots\\
a_0+a_1T_n+a_2T_n^2+\cdots+a_{n-1}T_n^{n-1}=d_n.
\end{cases}
\end{equation}

Ancora una volta, un problema apparentemente differente può essere
ricondotto alla soluzione di un sistema di equazioni lineari.


\subsection{La matrice di Vandermonde}

Il sistema dell'interpolazione polinomiale può essere scritto in
forma matriciale:

\begin{equation}
\underbrace{
\begin{pmatrix}
1&T_1&T_1^2&\cdots&T_1^{n-1}\\
1&T_2&T_2^2&\cdots&T_2^{n-1}\\
\vdots&\vdots&\vdots&&\vdots\\
1&T_n&T_n^2&\cdots&T_n^{n-1}
\end{pmatrix}
}_{A}
\underbrace{
\begin{pmatrix}
a_0\\
a_1\\
\vdots\\
a_{n-1}
\end{pmatrix}
}_{x}
=
\underbrace{
\begin{pmatrix}
d_1\\
d_2\\
\vdots\\
d_n
\end{pmatrix}
}_{b}.
\label{eq:vandermonde-system}
\end{equation}

La matrice

\begin{equation}
A=
\begin{pmatrix}
1&T_1&T_1^2&\cdots&T_1^{n-1}\\
1&T_2&T_2^2&\cdots&T_2^{n-1}\\
\vdots&\vdots&\vdots&&\vdots\\
1&T_n&T_n^2&\cdots&T_n^{n-1}
\end{pmatrix}
\end{equation}

prende il nome di \emph{matrice di Vandermonde}.

Il problema dell'interpolazione assume quindi nuovamente la forma

\begin{equation}
    Ax=b.
\end{equation}

Questa formulazione mostra ancora una volta la generalità degli
strumenti numerici studiati nel corso: la stessa struttura matematica
compare nella soluzione di problemi economici, ambientali e
finanziari e, ora, anche nella costruzione di una funzione interpolante.


\subsection{Perché non utilizzare semplicemente molti punti?}

A prima vista sembrerebbe naturale utilizzare il maggior numero
possibile di osservazioni.

Se disponiamo di $n$ BOT possiamo costruire un polinomio di grado
$n-1$ che passa esattamente attraverso tutti gli $n$ punti osservati.

Tuttavia, aumentando il numero di punti aumenta anche il grado del
polinomio e la matrice di Vandermonde contiene potenze sempre più
elevate delle maturities:

\begin{equation}
    1,\quad T,\quad T^2,\quad T^3,\quad\ldots,\quad T^{n-1}.
\end{equation}

Questo può generare un problema numerico importante: la matrice può
diventare \emph{mal condizionata}.

Ciò significa che piccole variazioni nei dati possono produrre
variazioni molto più grandi nella soluzione del sistema.

Questo aspetto è particolarmente rilevante quando lavoriamo con dati
reali. Prezzi, rendimenti e fattori di sconto sono infatti quantità
misurate numericamente e possono essere soggetti ad arrotondamenti,
errori di misura o piccole variazioni di mercato.


\subsection{Il numero di condizionamento}

Consideriamo nuovamente un sistema lineare

\begin{equation}
    Ax=b.
\end{equation}

Una misura della sensibilità numerica del problema è fornita dal
\emph{numero di condizionamento} della matrice $A$, definito come

\begin{equation}
    \boxed{
    \kappa(A)=\|A\|\,\|A^{-1}\|
    }.
    \label{eq:condition-number}
\end{equation}

Il numero di condizionamento ci dice quanto la soluzione del sistema
può essere sensibile a piccole perturbazioni nei dati.

Se il vettore dei termini noti subisce una perturbazione

\begin{equation}
    b\longrightarrow b+\Delta b,
\end{equation}

anche la soluzione cambia:

\begin{equation}
    x\longrightarrow x+\Delta x.
\end{equation}

Una relazione importante è

\begin{equation}
    \frac{\|\Delta x\|}{\|x\|}
    \leq
    \kappa(A)
    \frac{\|\Delta b\|}{\|b\|}.
    \label{eq:condition-bound}
\end{equation}

Questa espressione permette di comprendere il significato economico
e numerico del numero di condizionamento.

Se $\kappa(A)$ è piccolo, una piccola variazione nei dati produce
generalmente una piccola variazione nella soluzione. Se invece
$\kappa(A)$ è molto grande, una perturbazione apparentemente
trascurabile nei dati può essere fortemente amplificata.


\subsection{Interpretazione del condizionamento}

In termini qualitativi possiamo distinguere tre situazioni:

\begin{equation}
\begin{array}{lll}
\kappa(A)\approx 1
&\Rightarrow&
\text{sistema ben condizionato},\\[2mm]
\kappa(A)\gg 1
&\Rightarrow&
\text{sistema mal condizionato},\\[2mm]
\kappa(A)=\infty
&\Rightarrow&
\text{matrice singolare}.
\end{array}
\end{equation}

È importante distinguere il problema dell'esistenza della soluzione
dal problema della sua stabilità numerica.

Una matrice può infatti essere invertibile e il sistema può avere una
soluzione unica, ma la soluzione può essere estremamente sensibile a
piccole variazioni nei dati.

Pertanto,

\begin{quote}
avere una soluzione matematica unica non significa necessariamente
avere un problema numericamente ben posto.
\end{quote}

Questa osservazione è particolarmente importante per l'interpolazione
polinomiale. Utilizzare un numero crescente di punti permette di
costruire un polinomio che passa esattamente attraverso tutte le
osservazioni, ma può contemporaneamente peggiorare le proprietà
numeriche del problema.

Nasce quindi una nuova domanda: è possibile utilizzare molti punti
della funzione di sconto senza ricorrere a un unico polinomio di
grado elevato?

La risposta conduce a una diversa tecnica di interpolazione, basata
sull'utilizzo di polinomi di basso grado definiti localmente su
intervalli successivi: le \emph{spline}.

\section{Interpolazione mediante spline cubiche}
\label{sec:spline}

Nella sezione precedente abbiamo visto che l'interpolazione polinomiale
permette di costruire una funzione che passa esattamente attraverso un
insieme di punti osservati. Se disponiamo di $n$ punti, possiamo infatti
costruire un polinomio di grado $n-1$ determinandone i coefficienti
attraverso la soluzione di un sistema lineare basato sulla matrice di
Vandermonde.

Questa procedura presenta tuttavia un limite importante. All'aumentare
del numero di punti aumenta anche il grado del polinomio e la matrice
di Vandermonde può diventare mal condizionata. Di conseguenza, piccole
perturbazioni nei dati possono produrre variazioni rilevanti nei
coefficienti del polinomio e quindi nella funzione interpolata.

Sorge allora una domanda naturale: possiamo utilizzare molti punti
osservati senza costruire un unico polinomio di grado elevato?

Una possibile risposta consiste nell'abbandonare l'idea di utilizzare
un unico polinomio sull'intero intervallo e costruire invece diversi
polinomi di basso grado, ciascuno valido soltanto localmente.

Questa è l'idea fondamentale alla base delle \emph{spline}.


\subsection{L'idea delle spline}

Supponiamo di disporre di $n$ punti osservati

\begin{equation}
    (T_1,d_1),(T_2,d_2),\ldots,(T_n,d_n),
\end{equation}

ordinati in modo che

\begin{equation}
    T_1<T_2<\cdots<T_n.
\end{equation}

Le maturities $T_1,\ldots,T_n$ vengono chiamate \emph{nodi}.

Nell'interpolazione polinomiale globale cercheremmo un unico
polinomio

\begin{equation}
    p_{n-1}(T)
\end{equation}

di grado $n-1$ che passi attraverso tutti gli $n$ punti.

Con una spline adottiamo invece una strategia differente. Dividiamo
l'intervallo complessivo nei sottointervalli

\begin{equation}
[T_1,T_2],\quad
[T_2,T_3],\quad
\ldots,\quad
[T_{n-1},T_n]
\end{equation}

e associamo a ciascun intervallo una funzione polinomiale distinta.

In altre parole, non utilizziamo più un'unica funzione globale, ma una
funzione definita a tratti:

\begin{equation}
S(T)=
\begin{cases}
S_1(T), & T_1\leq T\leq T_2,\\
S_2(T), & T_2\leq T\leq T_3,\\
\vdots & \\
S_{n-1}(T), & T_{n-1}\leq T\leq T_n.
\end{cases}
\end{equation}

Il problema consiste quindi nel determinare i diversi polinomi
$S_i(T)$ e nel collegarli in modo tale che, nel loro insieme,
costituiscano una curva sufficientemente regolare.


\subsection{La spline cubica}

Una scelta molto comune consiste nell'utilizzare polinomi di terzo
grado. In questo caso si parla di \emph{spline cubica}.

Per ciascun intervallo $[T_i,T_{i+1}]$ definiamo

\begin{equation}
    S_i(T)
    =
    a_i+b_iT+c_iT^2+e_iT^3,
    \qquad i=1,\ldots,n-1.
    \label{eq:cubic-spline}
\end{equation}

Ogni intervallo possiede quindi quattro coefficienti:

\begin{equation}
    a_i,\quad b_i,\quad c_i,\quad e_i.
\end{equation}

Poiché $n$ nodi individuano $n-1$ intervalli, il numero complessivo
dei coefficienti da determinare è

\begin{equation}
    4(n-1).
    \label{eq:spline-number-coefficients}
\end{equation}

Nel nostro esempio finanziario disponiamo di 15 BOT. Abbiamo quindi

\begin{equation}
    n=15
\end{equation}

nodi e

\begin{equation}
    n-1=14
\end{equation}

intervalli.

Il numero totale dei coefficienti incogniti è pertanto

\begin{equation}
    4(15-1)=56.
\end{equation}

La costruzione diretta della spline richiede dunque la determinazione
di 56 coefficienti.

Il problema può sembrare notevolmente più complesso
dell'interpolazione lineare. Vedremo però che le proprietà che
desideriamo imporre alla spline generano precisamente le equazioni
necessarie per determinare questi coefficienti.


\subsection{Prima condizione: interpolazione dei nodi}

La prima proprietà è la stessa che caratterizza qualsiasi funzione
interpolante: la spline deve passare attraverso i punti osservati.

Consideriamo il polinomio $S_i(T)$ definito nell'intervallo

\begin{equation}
    [T_i,T_{i+1}].
\end{equation}

Esso deve assumere il valore $d_i$ nel primo estremo e il valore
$d_{i+1}$ nel secondo estremo. Dobbiamo quindi imporre

\begin{equation}
    S_i(T_i)=d_i
\end{equation}

e

\begin{equation}
    S_i(T_{i+1})=d_{i+1},
\end{equation}

per

\begin{equation}
    i=1,\ldots,n-1.
\end{equation}

Ogni intervallo fornisce quindi due equazioni. Poiché gli intervalli
sono $n-1$, otteniamo complessivamente

\begin{equation}
    2(n-1)
\end{equation}

equazioni.

Nel caso dei 15 BOT:

\begin{equation}
    2(15-1)=28.
\end{equation}

Le condizioni di interpolazione forniscono dunque 28 delle 56
equazioni necessarie.


\subsection{Seconda condizione: continuità della derivata prima}

Il semplice passaggio attraverso i nodi non è sufficiente.

Potremmo infatti costruire polinomi che si incontrano nei nodi ma che
presentano bruschi cambiamenti di pendenza. La funzione risultante
sarebbe continua, ma mostrerebbe degli ``spigoli'' nei punti di
raccordo.

Per ottenere una curva regolare richiediamo quindi che, in ciascun
nodo interno, i due polinomi adiacenti abbiano la stessa derivata
prima.

Nel nodo $T_i$ imponiamo pertanto

\begin{equation}
    S'_{i-1}(T_i)=S'_i(T_i),
    \qquad i=2,\ldots,n-1.
    \label{eq:spline-first-derivative}
\end{equation}

Partendo da

\begin{equation}
    S_i(T)=a_i+b_iT+c_iT^2+e_iT^3,
\end{equation}

la derivata prima è

\begin{equation}
    S'_i(T)
    =
    b_i+2c_iT+3e_iT^2.
\end{equation}

Le condizioni di continuità della derivata prima sono quindi
equazioni lineari nei coefficienti della spline.

I nodi interni sono $n-2$ e otteniamo pertanto altre

\begin{equation}
    n-2
\end{equation}

equazioni.

Con 15 nodi:

\begin{equation}
    15-2=13.
\end{equation}

Abbiamo quindi altre 13 equazioni.


\subsection{Terza condizione: continuità della derivata seconda}

Richiediamo una condizione ulteriore di regolarità.

Non vogliamo soltanto che i polinomi abbiano la stessa pendenza nei
punti nei quali si incontrano, ma anche che la curvatura della
funzione cambi in maniera regolare.

Imponiamo quindi la continuità della derivata seconda:

\begin{equation}
    S''_{i-1}(T_i)=S''_i(T_i),
    \qquad i=2,\ldots,n-1.
    \label{eq:spline-second-derivative}
\end{equation}

Poiché

\begin{equation}
    S'_i(T)=b_i+2c_iT+3e_iT^2,
\end{equation}

abbiamo

\begin{equation}
    S''_i(T)=2c_i+6e_iT.
\end{equation}

Anche queste condizioni sono lineari rispetto ai coefficienti
incogniti.

Poiché abbiamo $n-2$ nodi interni, otteniamo altre

\begin{equation}
    n-2
\end{equation}

equazioni.

Nel nostro esempio:

\begin{equation}
    15-2=13.
\end{equation}

Abbiamo quindi ottenuto altre 13 equazioni.


\subsection{Il conteggio delle equazioni}

A questo punto è utile verificare quante equazioni abbiamo costruito.

Le condizioni di interpolazione dei nodi forniscono

\begin{equation}
    2(n-1)
\end{equation}

equazioni.

La continuità della derivata prima fornisce

\begin{equation}
    n-2
\end{equation}

equazioni, mentre la continuità della derivata seconda ne fornisce
altre

\begin{equation}
    n-2.
\end{equation}

Complessivamente abbiamo quindi

\begin{equation}
\begin{split}
2(n-1)+(n-2)+(n-2)
&=2n-2+n-2+n-2\\
&=4n-6.
\end{split}
\end{equation}

Le incognite sono però

\begin{equation}
    4(n-1)=4n-4.
\end{equation}

Abbiamo dunque due incognite in più rispetto al numero di equazioni:

\begin{equation}
    (4n-4)-(4n-6)=2.
\end{equation}

Per determinare univocamente tutti i coefficienti della spline sono
quindi necessarie altre due condizioni.


\subsection{Condizioni al contorno e spline cubica naturale}

Le ultime due equazioni vengono ottenute imponendo delle condizioni
agli estremi dell'intervallo.

Una possibile scelta è quella della \emph{natural cubic spline}, o
spline cubica naturale.

In questo caso imponiamo che la derivata seconda sia nulla nei due
estremi:

\begin{equation}
    S''_1(T_1)=0
\end{equation}

e

\begin{equation}
    S''_{n-1}(T_n)=0.
\end{equation}

Queste due condizioni completano il sistema.

Abbiamo infatti

\begin{equation}
    4n-6+2=4n-4
\end{equation}

equazioni, esattamente pari al numero

\begin{equation}
    4(n-1)=4n-4
\end{equation}

di coefficienti incogniti.


\subsection{La spline come sistema lineare}

Possiamo ora tornare al nostro esempio con 15 BOT.

Il numero totale delle incognite è

\begin{equation}
    4(15-1)=56.
\end{equation}

Le diverse condizioni producono:

\begin{center}
\begin{tabular}{lr}
\hline
Condizione & Numero di equazioni\\
\hline
Interpolazione dei nodi & 28\\
Continuità di $S'$ & 13\\
Continuità di $S''$ & 13\\
Condizioni al contorno & 2\\
\hline
Totale & 56\\
\hline
\end{tabular}
\end{center}

Abbiamo quindi

\begin{equation}
    \boxed{
    56\text{ equazioni}
    \qquad
    56\text{ incognite}
    }.
\end{equation}

Se raccogliamo tutti i coefficienti dei polinomi in un unico vettore
$x$, tutte le condizioni precedenti possono essere organizzate in
forma matriciale.

Il problema assume quindi la forma

\begin{equation}
    \boxed{
    A_{56\times56}x_{56\times1}
    =
    b_{56\times1}
    }.
\end{equation}

Anche una procedura apparentemente complessa come la costruzione di
una spline cubica viene quindi ricondotta allo stesso problema
computazionale incontrato più volte nel corso:

\begin{equation}
    Ax=b.
\end{equation}

Questo risultato è particolarmente importante dal punto di vista
didattico. Il problema finanziario ci ha condotto alla necessità di
interpolare una funzione di sconto; la necessità di utilizzare molti
punti ci ha spinto a considerare le spline; la costruzione della
spline, a sua volta, si traduce nuovamente nella soluzione di un
sistema di equazioni lineari.

Le tecniche numeriche studiate in precedenza non costituiscono quindi
strumenti isolati, ma possono essere utilizzate come componenti di
procedure computazionali più articolate.

\section{Costruzione della struttura a termine con i BTP}
\label{sec:term-structure-btp}

Nelle sezioni precedenti abbiamo introdotto gli strumenti necessari
per affrontare il problema della costruzione della struttura a termine
dei tassi di interesse.

I BOT ci hanno permesso di osservare direttamente alcuni punti della
funzione di sconto. Abbiamo poi introdotto l'interpolazione per
ottenere valori della funzione in corrispondenza di scadenze non
direttamente osservate. In particolare, le spline cubiche permettono
di rappresentare la funzione di sconto mediante polinomi di basso
grado definiti localmente.

Possiamo ora tornare ai BTP e utilizzare questi strumenti per
costruire un'applicazione più vicina ai dati finanziari reali.

L'obiettivo è utilizzare congiuntamente:

\begin{itemize}
    \item i prezzi osservati dei BTP;
    \item i flussi di cassa generati dai titoli;
    \item la rappresentazione mediante spline della funzione di sconto;
\end{itemize}

per ricavare informazioni sulla struttura a termine.


\subsection{L'equazione di pricing di un BTP}

Consideriamo un BTP che genera $m$ pagamenti futuri alle date

\begin{equation}
    T_1,T_2,\ldots,T_m.
\end{equation}

Indicando con $CF_j$ il flusso di cassa alla data $T_j$, il prezzo
del titolo può essere scritto come

\begin{equation}
    P
    =
    \sum_{j=1}^{m}
    CF_jd(T_j),
    \label{eq:btp-pricing-general}
\end{equation}

dove $d(T_j)$ rappresenta il fattore di sconto corrispondente alla
data del pagamento.

Per un BTP nominale a tasso fisso, i pagamenti intermedi sono
costituiti dalle cedole, mentre l'ultimo pagamento comprende anche
il rimborso del valore nominale.

Se indichiamo con $C$ la cedola e assumiamo un valore nominale pari
a 100, possiamo quindi scrivere

\begin{equation}
    P
    =
    Cd(T_1)
    +Cd(T_2)
    +\cdots+
    (100+C)d(T_m).
    \label{eq:btp-pricing-fixed}
\end{equation}

Il problema è che i valori

\begin{equation}
    d(T_1),d(T_2),\ldots,d(T_m)
\end{equation}

non sono direttamente osservabili per tutte le date di pagamento.


\subsection{Rappresentare la funzione di sconto mediante una spline}

Nella sezione precedente abbiamo costruito una spline cubica per
rappresentare la funzione di sconto.

Supponiamo che la data $T_j$ appartenga all'intervallo compreso tra
due nodi consecutivi:

\begin{equation}
    T_i\leq T_j\leq T_{i+1}.
\end{equation}

All'interno di questo intervallo il fattore di sconto viene
rappresentato dal polinomio

\begin{equation}
    d(T)
    =
    a_i+b_iT+c_iT^2+e_iT^3.
    \label{eq:discount-spline-segment}
\end{equation}

Possiamo quindi calcolare il fattore di sconto relativo alla data
$T_j$ come

\begin{equation}
    d(T_j)
    =
    a_i+b_iT_j+c_iT_j^2+e_iT_j^3.
    \label{eq:discount-at-cashflow}
\end{equation}

Una volta conosciuti i coefficienti della spline, possiamo pertanto
ottenere il fattore di sconto corrispondente a qualsiasi data
compresa nell'intervallo coperto dai dati.


\subsection{Dalla spline all'equazione di pricing}

Consideriamo nuovamente l'equazione

\begin{equation}
    P
    =
    \sum_{j=1}^{m}
    CF_jd(T_j).
\end{equation}

Per ciascuna data $T_j$, il valore $d(T_j)$ è determinato dal tratto
della spline che contiene quella data.

Se, ad esempio, $T_j$ appartiene all'intervallo
$[T_i,T_{i+1}]$, possiamo sostituire

\begin{equation}
    d(T_j)
    =
    a_i+b_iT_j+c_iT_j^2+e_iT_j^3.
\end{equation}

Il contributo del flusso di cassa $CF_j$ al prezzo diventa quindi

\begin{equation}
    CF_jd(T_j)
    =
    CF_j
    \left(
    a_i+b_iT_j+c_iT_j^2+e_iT_j^3
    \right).
\end{equation}

Sviluppando il prodotto otteniamo

\begin{equation}
    CF_jd(T_j)
    =
    CF_ja_i
    +
    CF_jT_jb_i
    +
    CF_jT_j^2c_i
    +
    CF_jT_j^3e_i.
    \label{eq:cf-spline-linear}
\end{equation}

Questa espressione mette in evidenza una proprietà fondamentale:
l'equazione è \emph{lineare nei coefficienti della spline}.

Le quantità

\begin{equation}
    CF_j,\qquad
    CF_jT_j,\qquad
    CF_jT_j^2,\qquad
    CF_jT_j^3
\end{equation}

sono infatti note. Le incognite sono invece

\begin{equation}
    a_i,\qquad b_i,\qquad c_i,\qquad e_i.
\end{equation}

Questo permette di trasformare l'equazione di pricing di ciascun
BTP in un'equazione lineare.


\subsection{Il contributo dei diversi intervalli}

Un BTP genera normalmente più pagamenti e le relative date possono
appartenere a intervalli differenti della spline.

Consideriamo schematicamente tre intervalli:

\begin{equation}
[T_1,T_2],
\qquad
[T_2,T_3],
\qquad
[T_3,T_4].
\end{equation}

A essi corrispondono tre polinomi:

\begin{align}
S_1(T)&=a_1+b_1T+c_1T^2+e_1T^3,\\
S_2(T)&=a_2+b_2T+c_2T^2+e_2T^3,\\
S_3(T)&=a_3+b_3T+c_3T^2+e_3T^3.
\end{align}

Se un pagamento avviene nel primo intervallo, utilizzeremo i
coefficienti di $S_1$. Se avviene nel secondo intervallo,
utilizzeremo quelli di $S_2$, e così via.

L'equazione di pricing del titolo può quindi coinvolgere
coefficienti appartenenti a diversi segmenti della spline.

Possiamo raccogliere tutti questi coefficienti in un unico vettore:

\begin{equation}
a=
\begin{pmatrix}
a_1\\
b_1\\
c_1\\
e_1\\
a_2\\
b_2\\
c_2\\
e_2\\
\vdots
\end{pmatrix}.
\label{eq:spline-coefficient-vector-finance}
\end{equation}

Per ogni BTP possiamo allora costruire un vettore riga che contiene
i coefficienti noti che moltiplicano gli elementi di $a$.


\subsection{Un'equazione lineare per ciascun BTP}

Indichiamo con $B_i$ il vettore riga costruito utilizzando i flussi
di cassa e le date di pagamento del BTP $i$.

L'equazione di pricing può essere scritta come

\begin{equation}
    B_i a=P_i,
    \label{eq:single-btp-linear}
\end{equation}

dove:

\begin{itemize}
    \item $P_i$ è il prezzo osservato del titolo;
    \item $B_i$ contiene informazioni sui cash flow e sulle
          rispettive date;
    \item $a$ contiene i coefficienti incogniti della
          rappresentazione della funzione di sconto.
\end{itemize}

Ogni titolo osservato fornisce quindi una nuova equazione lineare.

Considerando simultaneamente più BTP possiamo raccogliere le
equazioni in un sistema:

\begin{equation}
\begin{pmatrix}
B_1\\
B_2\\
\vdots\\
B_N
\end{pmatrix}
a
=
\begin{pmatrix}
P_1\\
P_2\\
\vdots\\
P_N
\end{pmatrix}.
\end{equation}

Definendo

\begin{equation}
B=
\begin{pmatrix}
B_1\\
B_2\\
\vdots\\
B_N
\end{pmatrix}
\end{equation}

e

\begin{equation}
P=
\begin{pmatrix}
P_1\\
P_2\\
\vdots\\
P_N
\end{pmatrix},
\end{equation}

otteniamo

\begin{equation}
    \boxed{
    Ba=P
    }.
    \label{eq:BaP}
\end{equation}

Ancora una volta, il problema finanziario è stato trasformato nella
forma generale

\begin{equation}
    Ax=b.
\end{equation}


\subsection{Interpretazione della matrice $B$}

È utile soffermarsi sul significato della matrice $B$.

Le righe della matrice corrispondono ai titoli utilizzati
nell'applicazione. Le colonne corrispondono invece ai coefficienti
della funzione utilizzata per rappresentare i fattori di sconto.

Gli elementi di $B$ vengono costruiti combinando:

\begin{equation}
\boxed{
\text{cash flow}
\quad+\quad
\text{date dei cash flow}
\quad+\quad
\text{forma della spline}.
}
\end{equation}

La matrice $B$ rappresenta quindi il collegamento tra le
caratteristiche osservabili dei titoli e i parametri della funzione
di sconto che vogliamo determinare.

Questo passaggio è importante perché mostra come la costruzione di
un sistema lineare non sia un'operazione puramente formale.
La matrice dei coefficienti deve essere costruita a partire dalla
struttura economica o finanziaria del problema.


\subsection{Dalla soluzione del sistema alla funzione di sconto}

Una volta costruiti $B$ e $P$, possiamo risolvere

\begin{equation}
    Ba=P
\end{equation}

utilizzando le tecniche numeriche studiate nelle sezioni precedenti.

La soluzione fornisce il vettore dei coefficienti

\begin{equation}
    a^*.
\end{equation}

Conoscendo questi coefficienti possiamo ricostruire i diversi
polinomi della spline e quindi la funzione

\begin{equation}
    d(T).
\end{equation}

Il procedimento può essere sintetizzato come

\begin{equation}
\boxed{
\text{prezzi e cash flow}
\longrightarrow
Ba=P
\longrightarrow
a^*
\longrightarrow
d(T).
}
\end{equation}

A differenza del semplice esempio iniziale, non stiamo quindi
stimando separatamente un fattore di sconto per ciascuna data.
Determiniamo invece i parametri di una funzione che permette di
calcolare il fattore di sconto per le diverse scadenze.


\subsection{Dalla funzione di sconto alla curva spot}

Una volta ricostruita la funzione $d(T)$ possiamo trasformare i
fattori di sconto nei corrispondenti tassi spot.

Ricordiamo che

\begin{equation}
    d(T)
    =
    \frac{1}{(1+r(T))^T}.
\end{equation}

Risolvendo rispetto al tasso otteniamo

\begin{equation}
    \boxed{
    r(T)
    =
    d(T)^{-1/T}-1
    }.
    \label{eq:spline-discount-to-spot}
\end{equation}

Possiamo quindi scegliere una griglia di maturities

\begin{equation}
    T_1^*,T_2^*,\ldots,T_K^*
\end{equation}

e, per ciascuna di esse:

\begin{enumerate}
    \item calcolare il fattore di sconto $d(T_k^*)$;
    \item trasformarlo nel corrispondente tasso spot
          $r(T_k^*)$.
\end{enumerate}

Le coppie

\begin{equation}
    \left(T_k^*,r(T_k^*)\right)
\end{equation}

permettono infine di rappresentare graficamente la struttura a
termine dei tassi di interesse.


\subsection{Il ruolo dei BOT e dei BTP}

Possiamo ora comprendere meglio il ruolo svolto dai due tipi di
titoli utilizzati nell'applicazione.

I BOT, essendo titoli zero-coupon, permettono di ricavare
direttamente alcuni punti della funzione di sconto:

\begin{equation}
    d(T)=\frac{P_T}{100}.
\end{equation}

I BTP forniscono invece informazioni aggiuntive attraverso
l'intera sequenza dei loro flussi di cassa:

\begin{equation}
    P_i
    =
    \sum_j CF_{ij}d(T_{ij}).
\end{equation}

La funzione di interpolazione permette di collegare queste
informazioni e di rappresentare $d(T)$ anche in corrispondenza
delle date per le quali non osserviamo direttamente un titolo
zero-coupon.

La costruzione della curva può quindi essere vista come
l'integrazione di tre elementi:

\begin{equation}
\boxed{
\begin{array}{ccc}
\text{BOT}
&
+
&
\text{BTP}\\[1mm]
\downarrow&&\downarrow\\
\text{punti della funzione di sconto}
&&
\text{equazioni di pricing}\\[2mm]
&\searrow\quad\swarrow&\\[-1mm]
&d(T)&
\end{array}}
\end{equation}


\subsection{Dal problema finanziario al problema computazionale}

L'intero percorso seguito nell'applicazione può essere ora
ricostruito.

Siamo partiti da una domanda finanziaria:

\begin{quote}
come possiamo ricavare la struttura a termine dei tassi di interesse
utilizzando i prezzi dei titoli osservati sul mercato?
\end{quote}

Abbiamo inizialmente considerato il caso semplice dei BOT, per i
quali il prezzo permette di ottenere direttamente un fattore di
sconto.

Successivamente abbiamo incontrato un problema: le scadenze
disponibili sul mercato non coincidono necessariamente con tutte le
date nelle quali abbiamo bisogno di conoscere il fattore di sconto.

Questo ci ha condotto all'interpolazione:

\begin{equation}
\text{punti osservati}
\longrightarrow
\text{funzione }d(T).
\end{equation}

L'interpolazione polinomiale globale ha evidenziato possibili
problemi di condizionamento. Abbiamo quindi introdotto le spline,
che rappresentano la funzione mediante polinomi di basso grado
definiti localmente.

Infine, siamo tornati alle equazioni di pricing dei BTP. Sostituendo
in tali equazioni la rappresentazione scelta per $d(T)$ abbiamo
ottenuto un sistema lineare:

\begin{equation}
    Ba=P.
\end{equation}

La soluzione del sistema fornisce i coefficienti necessari per
ricostruire la funzione di sconto e, successivamente, la curva dei
tassi spot.


\subsection{Il messaggio computazionale}

L'applicazione finanziaria può essere sintetizzata attraverso la
sequenza

\begin{equation}
\boxed{
\text{prezzi + cash flow}
\longrightarrow
Ba=P
\longrightarrow
a
\longrightarrow
d(T)
\longrightarrow
r(T).
}
\end{equation}

Il passaggio centrale è ancora una volta la soluzione di un sistema
lineare.

Tuttavia, questa applicazione mette in evidenza anche un secondo
aspetto importante del calcolo numerico: la qualità della soluzione
non dipende soltanto dall'algoritmo utilizzato per risolvere il
sistema.

Dipende anche da \emph{come il problema viene rappresentato}.

Nel nostro caso, la scelta della funzione utilizzata per descrivere
$d(T)$ determina la struttura del sistema lineare che dobbiamo
risolvere e influenza le proprietà della curva ottenuta.

L'interpolazione polinomiale e le spline non sono quindi semplici
passaggi accessori, ma fanno parte della formulazione numerica del
problema.

L'applicazione finanziaria riunisce così i principali concetti
introdotti in questa parte del corso:

\begin{equation}
\boxed{
\begin{array}{c}
\text{problema finanziario}\\
\downarrow\\
\text{rappresentazione matematica}\\
\downarrow\\
\text{interpolazione}\\
\downarrow\\
\text{sistema lineare}\\
\downarrow\\
\text{soluzione numerica}\\
\downarrow\\
\text{interpretazione finanziaria}
\end{array}}
\end{equation}

Il risultato finale non è quindi soltanto una curva dei tassi.
L'esempio mostra, più in generale, come un problema finanziario reale
possa essere trasformato progressivamente in una sequenza di problemi
matematici che il computer è in grado di risolvere.

\end{document}
